数组、双指针和滑动窗口深度题卡 按题型拆成章内大节，但每个大节都先从 LeetCode 案例切入，再进入分流、案例矩阵、案例推导和实验复盘。这样读者看到的是从具体题推到规律，而不是先读抽象方法再猜它能用在哪。

\topic{数组与原地维护}
这一节先看 LeetCode 5 最长回文子串、LeetCode 26 删除有序数组重复项、LeetCode 27 移除元素 这几道 LeetCode 题，再整理同一题型的分流和推导。阅读顺序不是先背原理，而是先看案例的暴力瓶颈，再把重复出现的观察抽象成方法。

\subsection{原理和适用条件}

以 LeetCode 5 最长回文子串、LeetCode 26 删除有序数组重复项、LeetCode 27 移除元素 为主线：LeetCode 5 最长回文子串：模型是“回文围绕中心对称，中心可以是字符也可以是两个字符之间。”，优化抓手是“中心扩展避免枚举子串后再逐个检查；每次扩展只比较新增左右边界。”；LeetCode 26 删除有序数组重复项：模型是“有序数组让重复元素相邻，慢指针表示下一个唯一值写入位置。”，优化抓手是“快指针扫描，只有发现新值才写入慢指针并推进。”；LeetCode 27 移除元素：模型是“原地过滤，慢指针左侧始终是不等于目标值的稳定前缀。”，优化抓手是“保持顺序用快慢指针；不保持顺序可用左右交换减少写入。”。这些案例共同推出的规律是：先写一个不会漏答案的枚举或辅助数组版本，再观察哪些位置已经被前缀、后缀、慢指针或边界确定。优化时把数组划成语义清楚的区域，每次循环只扩大已确认区域。 方法名只是复盘后的标签，真正要掌握的是从题面约束、暴力失败、状态设计到边界反例的推导链。

\begin{principlebox}{图示：从 LeetCode 案例到 数组与原地维护推导链}
\begin{tabular}{@{}p{0.18\linewidth}p{0.74\linewidth}@{}}
暴力基线 & 枚举候选、完整模拟或逐个检查，先保证覆盖全部可能答案。\\
瓶颈定位 & 慢点在于重复搬移和重复求同一段信息。只要一个位置的状态已经由前缀、后缀、慢指针或边界指针确定，再回头扫描就是浪费。\\
结构选择 & 优化要把数组切成若干语义明确的区间：已经处理好的前缀、未知区、已经确定的后缀，或者左侧贡献和右侧贡献。双指针、前后缀数组、原地交换和稳定写入，本质都是让每个元素只进入少数几次状态转移。\\
不变量 & 不变量必须能落到下标上：慢指针左侧保存什么，右指针右侧保存什么，当前下标之前是否已经满足题目要求。只要这个叙述含糊，代码中的加一减一就会变成猜。\\
代码落地 & C++ 实现优先使用 \code{std::vector<int>} 和明确的整型下标。涉及乘积、面积、总和时要主动判断是否需要 \code{long long}。如果题目要求原地修改，要说明返回值表示新长度还是数组本身，避免把输出语义和容器容量混在一起。\\
\end{tabular}
\end{principlebox}

\subsection{LeetCode 案例分流}

\begin{principlebox}{从 LeetCode 案例分流：数组与原地维护}
\textbf{原地过滤和稳定写入。}
代表案例：LeetCode 27 移除元素、LeetCode 26 删除有序数组重复项、LeetCode 283 移动零。先看这些题的题面条件：题目要求在数组本身上删除、移动或压缩元素，并且通常只返回有效长度。由此推出的处理方式是：用慢指针表示下一个写入位置，快指针扫描所有输入元素。风险点：必须区分稳定顺序和不稳定顺序，返回长度不等于容器容量。

\textbf{前后缀贡献。}
代表案例：LeetCode 238 除自身以外数组的乘积、LeetCode 42 接雨水、LeetCode 581 最短无序连续子数组。先看这些题的题面条件：当前位置答案由左侧信息和右侧信息共同决定，且两侧可以独立累积。由此推出的处理方式是：先保存一侧贡献，再反向扫描合并另一侧贡献。风险点：零、负数、溢出和边界位置会改变初始化语义。

\textbf{排列和局部字典序。}
代表案例：LeetCode 31 下一个排列、LeetCode 54 螺旋矩阵、LeetCode 88 合并两个有序数组。先看这些题的题面条件：题目要求找下一个排列、局部重排或保持全局字典序最小变化。由此推出的处理方式是：从后向前找可调整的枢轴，再让后缀恢复最小结构。风险点：后缀处理顺序错了会得到一个更大的但不是下一个的答案。

\textbf{回文和中心结构。}
代表案例：LeetCode 5 最长回文子串、LeetCode 647 回文子串。先看这些题的题面条件：字符串或数组片段具有左右对称关系。由此推出的处理方式是：枚举中心或中心间隙，向两侧扩展，只比较新增边界。风险点：回文子串和回文子序列不是同一问题，后者通常转为区间 DP。

\end{principlebox}

\subsection{案例矩阵}

本节案例覆盖 LeetCode 5 最长回文子串、LeetCode 26 删除有序数组重复项、LeetCode 27 移除元素、LeetCode 28 找出字符串中第一个匹配项的下标、LeetCode 31 下一个排列、LeetCode 41 缺失的第一个正数、LeetCode 54 螺旋矩阵、LeetCode 75 颜色分类、LeetCode 88 合并两个有序数组、LeetCode 189 轮转数组、LeetCode 204 计数质数、LeetCode 238 除自身以外数组的乘积、LeetCode 283 移动零、LeetCode 304 二维区域和检索 - 矩阵不可变、LeetCode 307 区域和检索 - 数组可修改、LeetCode 344 反转字符串、LeetCode 459 重复的子字符串、LeetCode 493 翻转对、LeetCode 581 最短无序连续子数组、LeetCode 587 安装栅栏、LeetCode 647 回文子串、LeetCode 1122 数组的相对排序。阅读时不要按题号背答案，而要先判断题面落在哪个分流场景：查询目标是什么，输入约束给了什么单调性或等价关系，暴力慢在哪里，优化结构保存了什么状态。

\begin{principlebox}{案例矩阵}
\textbf{LeetCode 5 最长回文子串。}
回文围绕中心对称，中心可以是字符也可以是两个字符之间。单点风险：偶数长度回文、单字符、全相同字符、多个同长答案。

\textbf{LeetCode 26 删除有序数组重复项。}
有序数组让重复元素相邻，慢指针表示下一个唯一值写入位置。单点风险：空数组、全重复、无重复、返回长度不等于物理容量。

\textbf{LeetCode 27 移除元素。}
原地过滤，慢指针左侧始终是不等于目标值的稳定前缀。单点风险：所有元素都被移除、没有元素被移除、目标值在尾部密集。

\textbf{LeetCode 28 找出字符串中第一个匹配项的下标。}
字符串匹配的暴力回退会反复比较已经匹配过的前缀。单点风险：空模式串、模式比文本长、重复字符、完全不匹配、重叠匹配。

\textbf{LeetCode 31 下一个排列。}
字典序结构由最长非递增后缀决定，枢轴是后缀前一个位置。单点风险：完全降序、重复数字、只有一个元素、交换后后缀必须最小化。

\textbf{LeetCode 41 缺失的第一个正数。}
答案只可能落在一到 n 加一之间，数组本身可以充当值到位置的映射。单点风险：非正数、大于 n 的数、重复值导致无限交换、数组已经完整包含一到 n。

\textbf{LeetCode 54 螺旋矩阵。}
四个边界收缩模拟一圈一圈访问。单点风险：单行、单列、空矩阵、最后一圈只剩一个元素。

\textbf{LeetCode 75 颜色分类。}
三指针维护零区、一号区、未知区和二区。单点风险：全零、全二、已经有序、逆序、交换后漏检查。

\textbf{LeetCode 88 合并两个有序数组。}
第一个数组尾部有空位，适合从后往前写入最终最大值。单点风险：第二个数组为空、第一个有效元素为空、重复值、尾部剩余。

\textbf{LeetCode 189 轮转数组。}
给定数组和整数 k，将数组中的元素整体向右轮转 k 个位置。单点风险：k 为零、k 大于 n、空数组、单元素。

\textbf{LeetCode 204 计数质数。}
逐个数试除会重复判断很多合数的因子。单点风险：n 小于二、从 i*i 开始防溢出、重复标记、边界是不包含 n。

\textbf{LeetCode 238 除自身以外数组的乘积。}
给定整数数组，返回 answer，其中 answer[i] 等于 nums[i] 之外所有元素的乘积。单点风险：一个零、多个零、负数、乘积溢出。

\textbf{LeetCode 283 移动零。}
给定数组，原地把所有 0 移到末尾，同时保持非零元素的相对顺序。单点风险：全零、无零、零在开头结尾、稳定顺序要求。

\textbf{LeetCode 304 二维区域和检索 - 矩阵不可变。}
固定矩阵上需要多次查询闭矩形区域和。单点风险：第一行第一列、单行单列、空矩阵、多次重复查询、坐标加一。

\textbf{LeetCode 307 区域和检索 - 数组可修改。}
单点更新会破坏普通前缀和后面所有位置。单点风险：下标一基转换、负数更新、重复更新、查询单点、差值更新。

\textbf{LeetCode 344 反转字符串。}
反转只需要把对称位置交换，没必要新建完整副本。单点风险：空数组、单字符、偶数长度、奇数长度、原地修改。

\textbf{LeetCode 459 重复的子字符串。}
判断字符串能否由某个较短周期重复组成，本质是找合法周期长度。单点风险：长度一、全相同字符、无重复、部分前后缀相同但不能整除。

\textbf{LeetCode 493 翻转对。}
每个 i 需要统计右侧满足 nums[i] 大于 2*nums[j] 的元素。单点风险：负数、乘以二溢出、重复值、严格大于、答案数量很大。

\textbf{LeetCode 581 最短无序连续子数组。}
要找破坏全局有序性的最左和最右位置。单点风险：已经有序、逆序、重复值、局部乱序。

\textbf{LeetCode 587 安装栅栏。}
外层点构成凸包，内部点不会影响围栏边界。单点风险：所有点共线、重复点、共线边界是否保留、叉积溢出。

\textbf{LeetCode 647 回文子串。}
每个回文都有唯一中心或中心间隙。单点风险：空串、全相同字符、偶数回文、奇数回文。

\textbf{LeetCode 1122 数组的相对排序。}
arr2 给出部分值的自定义顺序，剩余值按自然升序。单点风险：arr2 元素互不相同、arr1 中有不在 arr2 的值、重复值、值域边界、空数组。

\end{principlebox}

\subsection{共性落点}

\begin{principlebox}{本组共性落点}
\textbf{慢版本。} 暴力做法通常枚举每个位置的最终去向，或者为每个位置重新扫描一段区间。它容易写出正确答案，却没有利用数组的连续存储、下标可随机访问、前后缀可复用这些条件。
\textbf{瓶颈。} 慢点在于重复搬移和重复求同一段信息。只要一个位置的状态已经由前缀、后缀、慢指针或边界指针确定，再回头扫描就是浪费。
\textbf{状态字段。} 当前下标、已处理前缀边界、未知区边界、答案变量；若有前后缀贡献，还要说明左侧累计值和右侧累计值分别何时更新。
\textbf{不变量。} 不变量必须能落到下标上：慢指针左侧保存什么，右指针右侧保存什么，当前下标之前是否已经满足题目要求。只要这个叙述含糊，代码中的加一减一就会变成猜。
\textbf{实现检查。} 先写清数组长度为 0 或 1 的入口；主循环只推进一个明确边界；答案更新位置要和已处理区间一致。
\textbf{C++ 落地。} C++ 实现优先使用 \code{std::vector<int>} 和明确的整型下标。涉及乘积、面积、总和时要主动判断是否需要 \code{long long}。如果题目要求原地修改，要说明返回值表示新长度还是数组本身，避免把输出语义和容器容量混在一起。
\textbf{证明抓手。} 证明从区间不变量开始：已处理区域已经满足题意，未知区域还没有承诺，每次推进不会破坏已处理区域。
\end{principlebox}

\subsection{案例推导}

\noindent\textbf{LeetCode 5 最长回文子串。}

\textbf{题目说明。} LeetCode 5 最长回文子串 的题面入口要先还原成具体任务：回文围绕中心对称，中心可以是字符也可以是两个字符之间。

\textbf{样例入口。} 拿 aba 和 abba 手算，回文不是任意子串都要重查，而是围绕中心向两边同步比较。aba 的中心是 b，abba 的中心在两个 b 中间；每扩一层只新增左右两个字符。

\textbf{暴力基线。} 暴力枚举所有子串 s[left..right]，再用双指针检查它是否回文，记录最长回文子串。

\textbf{暴力过程。} 以 s=babad 为例，暴力会检查 b、ba、bab、baba、babad，再从 a 开始检查 a、ab、aba。bab 和 aba 的内部字符被反复比较。

\textbf{重复工作。} 题面模型：回文围绕中心对称，中心可以是字符也可以是两个字符之间。暴力版的慢点要落到本题：浪费在每个子串都从头判断回文，明明回文的合法性只取决于左右边界是否相等，以及内部是否已经是回文。后面的优化只消掉这类重复工作，不能改变答案集合。

\textbf{优化条件。} 本题的关键观察是：中心扩展避免枚举子串后再逐个检查；每次扩展只比较新增左右边界。这句话必须能翻译成可维护状态；如果题目条件改变后观察不再成立，就不能继续套原来的写法。

\textbf{相邻状态。} 规律是枚举 2n-1 个中心，每个中心只在新增边界相等时继续扩展。把这个特例继续拆开看：先标出已经确认的前缀或后缀，再标出未知区；能被丢掉的候选，必须是以后再也不会改变已确认区域语义的候选。这样才算从例子推出数组不变量，而不是只记一次操作。

\textbf{状态复用。} 把子串枚举改成中心枚举。中心可以是一个字符，也可以是两个字符之间；每次只向左右各扩一格，成功就更新答案边界。手算时只改被新输入影响的那一小部分，而不是回头重算完整候选。

\textbf{指针和字段。} 状态字段要能说清：center 表示回文中心，left/right 是当前扩展边界，best\_left/best\_len 记录最长回文的位置和长度。手算时按这个协议跟踪：aba 从中心 b 开始，先看 b，再扩到 a..a 成功；abba 从两个 b 中间开始，先看 bb，再扩到 a..a 成功。

\textbf{代码落点。} 关键代码是对每个 i 调用 expand(i, i) 和 expand(i, i+1)。expand 里 while left>=0 \&\& right<n \&\& s[left]==s[right]，再更新最长长度。关键代码行必须能逐句翻译成“旧状态怎样变成新状态”，否则就只是模板。

\textbf{正确性。} 任意回文都有唯一中心或中心间隙；从这个中心向外扩展时会经过它的全部边界，所以不会漏掉任何回文子串。

\textbf{边界实验。} 先用偶数长度回文、单字符、全相同字符、多个同长答案做反例检查。实验时把偶数长度回文、单字符、全相同字符、多个同长答案列成表，再打印关键下标和有效区间。若某个元素被写入后又被未来步骤破坏，说明区间不变量没有成立。

\textbf{复杂度变化。} 暴力枚举加检查可到 O(\ensuremath{n^{3}})，中心扩展有 O(n) 个中心、每个中心最多扩 O(n)，时间 O(\ensuremath{n^{2}})，额外空间 O(1)。

\textbf{迁移追问。} 如果题目改成“若要求回文子序列，应转成区间 DP，不能用中心扩展”，先判断原状态是否仍然覆盖新答案集合，再决定是微调边界、改 value 语义，还是换算法。

\noindent\textbf{LeetCode 26 删除有序数组重复项。}

\textbf{题目说明。} LeetCode 26 删除有序数组重复项 的题面入口要先还原成具体任务：有序数组让重复元素相邻，慢指针表示下一个唯一值写入位置。

\textbf{样例入口。} 拿 [1, 1, 2] 手算，write 先停在第一个 1 后面；读到第二个 1 时，它和已保留尾值相同，不写入；读到 2 时写到 write 位置，前缀变成 [1, 2]。

\textbf{暴力基线。} 暴力可以新建数组，把每组相同元素只写一次；或者每发现重复就移动后续元素覆盖它。

\textbf{暴力过程。} 以 [1, 1, 2] 为例，辅助数组写入 1，第二个 1 跳过，遇到 2 写入。原地做法只是把辅助数组的写入位置改成 nums[write]。

\textbf{重复工作。} 题面模型：有序数组让重复元素相邻，慢指针表示下一个唯一值写入位置。暴力版的慢点要落到本题：浪费在额外数组或反复搬移。数组已经有序，重复元素相邻，只需要和已保留前缀的最后一个值比较。后面的优化只消掉这类重复工作，不能改变答案集合。

\textbf{优化条件。} 本题的关键观察是：快指针扫描，只有发现新值才写入慢指针并推进。这句话必须能翻译成可维护状态；如果题目条件改变后观察不再成立，就不能继续套原来的写法。

\textbf{相邻状态。} 规律是慢指针左侧始终是不重复前缀，快指针只在遇到新值时推进慢指针；空数组和全重复数组只改变初始化，不改变不变量。把这个特例继续拆开看：先标出已经确认的前缀或后缀，再标出未知区；能被丢掉的候选，必须是以后再也不会改变已确认区域语义的候选。这样才算从例子推出数组不变量，而不是只记一次操作。

\textbf{状态复用。} write 表示下一个唯一值写入位置，read 从左到右扫描。nums[0..write-1] 始终是不重复前缀。手算时只改被新输入影响的那一小部分，而不是回头重算完整候选。

\textbf{指针和字段。} 状态字段要能说清：read 是当前读入位置，write 是有效前缀长度，nums[write-1] 是当前已保留的最后一个唯一值。手算时按这个协议跟踪：[1, 1, 2] 中 write 初始为 1；read=1 时 nums[1]==nums[0] 不写；read=2 时不同，写到 nums[1]，write 变 2。

\textbf{代码落点。} 关键代码是 if (nums[read] != nums[write-1]) nums[write++] = nums[read]; 。空数组直接返回 0。关键代码行必须能逐句翻译成“旧状态怎样变成新状态”，否则就只是模板。

\textbf{正确性。} 因为数组有序，同一个值只会连续出现；第一次出现被写入后，后续相同值都应跳过，遇到新值就扩展不重复前缀。

\textbf{边界实验。} 先用空数组、全重复、无重复、返回长度不等于物理容量做反例检查。实验时把空数组、全重复、无重复、返回长度不等于物理容量列成表，再打印关键下标和有效区间。若某个元素被写入后又被未来步骤破坏，说明区间不变量没有成立。

\textbf{复杂度变化。} 移动元素暴力最坏 O(\ensuremath{n^{2}}) 或辅助数组 O(n) 空间；快慢指针 O(n) 时间、O(1) 额外空间。

\textbf{迁移追问。} 如果题目改成“若允许每个元素最多保留两次，只需改变写入条件为和前两个比较”，先判断原状态是否仍然覆盖新答案集合，再决定是微调边界、改 value 语义，还是换算法。

\noindent\textbf{LeetCode 27 移除元素。}

\textbf{题目说明。} LeetCode 27 移除元素 的题面入口要先还原成具体任务：原地过滤，慢指针左侧始终是不等于目标值的稳定前缀。

\textbf{样例入口。} 拿 [3, 2, 2, 3]、val=3 手算，稳定写法读到 3 不写，读到两个 2 依次写入前缀，最后有效区间是 [2, 2]。若题目不要求顺序，也可以把左侧 3 和右侧非 3 交换来减少写入。

\textbf{暴力基线。} 暴力可以新建数组，只复制不等于 val 的元素；最后返回新数组长度或把结果写回原数组前缀。

\textbf{暴力过程。} 以 nums=[3, 2, 2, 3]、val=3 为例，辅助数组依次跳过 3、写入 2、写入 2、跳过 3，得到有效前缀 [2, 2]。

\textbf{重复工作。} 题面模型：原地过滤，慢指针左侧始终是不等于目标值的稳定前缀。暴力版的慢点要落到本题：浪费在辅助数组。若只需要原地返回有效长度，辅助数组的写入位置就是原数组里的 write 指针。后面的优化只消掉这类重复工作，不能改变答案集合。

\textbf{优化条件。} 本题的关键观察是：保持顺序用快慢指针；不保持顺序可用左右交换减少写入。这句话必须能翻译成可维护状态；如果题目条件改变后观察不再成立，就不能继续套原来的写法。

\textbf{相邻状态。} 规律是先确认是否稳定，再决定快慢指针还是左右交换；所有元素都被移除时返回长度 0。把这个特例继续拆开看：先标出已经确认的前缀或后缀，再标出未知区；能被丢掉的候选，必须是以后再也不会改变已确认区域语义的候选。这样才算从例子推出数组不变量，而不是只记一次操作。

\textbf{状态复用。} 稳定写法维护 write，nums[0..write-1] 都是不等于 val 的元素且相对顺序不变；read 扫描所有元素。手算时只改被新输入影响的那一小部分，而不是回头重算完整候选。

\textbf{指针和字段。} 状态字段要能说清：read 是输入扫描位置，write 是下一个保留元素写入位置，返回值 write 是有效长度，不是数组物理容量。手算时按这个协议跟踪：读到第一个 3 不写；读到 2 写到 nums[0]；读到第二个 2 写到 nums[1]；最后 3 不写，返回 2。

\textbf{代码落点。} 关键代码是 if (nums[read] != val) nums[write++] = nums[read]; 。若题目明确不要求稳定顺序，才考虑左右交换版本。关键代码行必须能逐句翻译成“旧状态怎样变成新状态”，否则就只是模板。

\textbf{正确性。} write 左侧只接收读过且不等于 val 的元素，read 每前进一步都保持这个前缀语义，扫描结束时前缀正好包含所有应保留元素。

\textbf{边界实验。} 先用所有元素都被移除、没有元素被移除、目标值在尾部密集做反例检查。实验时把所有元素都被移除、没有元素被移除、目标值在尾部密集列成表，再打印关键下标和有效区间。若某个元素被写入后又被未来步骤破坏，说明区间不变量没有成立。

\textbf{复杂度变化。} 辅助数组 O(n) 时间和 O(n) 空间；原地快慢指针 O(n) 时间、O(1) 额外空间。

\textbf{迁移追问。} 如果题目改成“工程里要先确认是否要求稳定顺序，不同要求对应不同写法”，先判断原状态是否仍然覆盖新答案集合，再决定是微调边界、改 value 语义，还是换算法。

\noindent\textbf{LeetCode 28 找出字符串中第一个匹配项的下标。}

\textbf{题目说明。} LeetCode 28 找出字符串中第一个匹配项的下标 的题面入口要先还原成具体任务：字符串匹配的暴力回退会反复比较已经匹配过的前缀。

\textbf{样例入口。} 拿 haystack=sadbutsad、needle=sad 手算，暴力从每个起点尝试匹配，起点 0 成功；若文本和模式有大量重复前缀，失配后从下一个起点重比会浪费。

\textbf{暴力基线。} 暴力从 haystack 的每个起点尝试匹配 needle，遇到失配就把起点加一，needle 从头再比。

\textbf{暴力过程。} 以 haystack=sadbutsad、needle=sad 为例，起点 0 成功。若换成很多重复前缀的文本，如 aaaaaab 匹配 aaaab，暴力会反复比较相同的 a。

\textbf{重复工作。} 题面模型：字符串匹配的暴力回退会反复比较已经匹配过的前缀。暴力版的慢点要落到本题：浪费在失配后完全丢掉已经匹配的前缀信息。模式串内部若有相同前后缀，下一次匹配可以复用一部分。后面的优化只消掉这类重复工作，不能改变答案集合。

\textbf{优化条件。} 本题的关键观察是：KMP 用前缀函数记录模式串最长相等前后缀，失配时把模式串跳到可复用前缀位置。这句话必须能翻译成可维护状态；如果题目条件改变后观察不再成立，就不能继续套原来的写法。

\textbf{相邻状态。} 规律是 KMP 用模式串前缀函数记录可复用的最长相等前后缀，失配时移动模式而不是回退文本。把这个特例继续拆开看：先标出已经确认的前缀或后缀，再标出未知区；能被丢掉的候选，必须是以后再也不会改变已确认区域语义的候选。这样才算从例子推出数组不变量，而不是只记一次操作。

\textbf{状态复用。} KMP 预处理 lps[i]，表示 needle[0..i] 的最长相等真前后缀长度。匹配时文本指针不回退，模式指针按 lps 跳转。手算时只改被新输入影响的那一小部分，而不是回头重算完整候选。

\textbf{指针和字段。} 状态字段要能说清：i 是文本位置，j 是模式位置，lps[j-1] 是失配后可复用的模式前缀长度。手算时按这个协议跟踪：模式 aaaab 在第 5 个字符失配时，前面 aaa 已经可作为下一轮前缀，不必让文本回到旧起点重新比较。

\textbf{代码落点。} 关键代码是失配时 if (j>0) j=lps[j-1]; else ++i。匹配到 j==m 时返回 i-m。关键代码行必须能逐句翻译成“旧状态怎样变成新状态”，否则就只是模板。

\textbf{正确性。} lps 保存的是已经匹配部分中最长可复用前缀，跳转后所有被跳过的起点都无法比这个前缀匹配更多字符。

\textbf{边界实验。} 先用空模式串、模式比文本长、重复字符、完全不匹配、重叠匹配做反例检查。实验时把空模式串、模式比文本长、重复字符、完全不匹配、重叠匹配列成表，再打印关键下标和有效区间。若某个元素被写入后又被未来步骤破坏，说明区间不变量没有成立。

\textbf{复杂度变化。} 暴力最坏 O(nm)，KMP 预处理 O(m)、匹配 O(n)，空间 O(m)。

\textbf{迁移追问。} 如果题目改成“若只做少量短串匹配，暴力可接受；大量匹配再考虑 KMP、哈希或自动机”，先判断原状态是否仍然覆盖新答案集合，再决定是微调边界、改 value 语义，还是换算法。

\noindent\textbf{LeetCode 31 下一个排列。}

\textbf{题目说明。} LeetCode 31 下一个排列 的题面入口要先还原成具体任务：字典序结构由最长非递增后缀决定，枢轴是后缀前一个位置。

\textbf{样例入口。} 拿 [1, 3, 2] 手算，后缀 [3, 2] 已经非递增，枢轴是 1；从右找刚好比 1 大的 2 交换，得到 [2, 3, 1]，再反转后缀成 [2, 1, 3]。

\textbf{暴力基线。} 暴力可以生成所有排列并排序，找到当前排列的下一个；这能定义答案，但复杂度完全不可接受。

\textbf{暴力过程。} 以 [1, 3, 2] 为例，所有更大的排列里最小的是 [2, 1, 3]。从右看，后缀 [3, 2] 已经是降序，无法在不改前缀的情况下变得更大。

\textbf{重复工作。} 题面模型：字典序结构由最长非递增后缀决定，枢轴是后缀前一个位置。暴力版的慢点要落到本题：浪费在枚举全部排列。真正需要改变的是最靠右、仍能被提高的一位；它右侧后缀已经是当前前缀下的最大排列。后面的优化只消掉这类重复工作，不能改变答案集合。

\textbf{优化条件。} 本题的关键观察是：从右找第一个上升点，再从右找刚好更大的元素交换，最后反转后缀。这句话必须能翻译成可维护状态；如果题目条件改变后观察不再成立，就不能继续套原来的写法。

\textbf{相邻状态。} 规律是最长非递增后缀已经是该前缀下的最大排列，必须提升枢轴并把后缀变成最小结构；完全降序时整体反转成最小排列。把这个特例继续拆开看：先标出已经确认的前缀或后缀，再标出未知区；能被丢掉的候选，必须是以后再也不会改变已确认区域语义的候选。这样才算从例子推出数组不变量，而不是只记一次操作。

\textbf{状态复用。} 从右向左找第一个 nums[i]<nums[i+1] 的枢轴 i；再从右找第一个大于 nums[i] 的数交换；最后反转后缀让它最小。手算时只改被新输入影响的那一小部分，而不是回头重算完整候选。

\textbf{指针和字段。} 状态字段要能说清：pivot 是要提高的位置，successor 是右侧刚好更大的交换对象，suffix 是交换后必须恢复升序的后缀。手算时按这个协议跟踪：[1, 3, 2] 中 pivot=1，successor=2，交换得 [2, 3, 1]，反转后缀 [3, 1] 得 [2, 1, 3]。

\textbf{代码落点。} 关键代码是先找 pivot；若不存在说明整个数组降序，直接 reverse 全数组。交换后 reverse(nums.begin()+pivot+1, nums.end())。关键代码行必须能逐句翻译成“旧状态怎样变成新状态”，否则就只是模板。

\textbf{正确性。} 最长非递增后缀已经是该前缀下的最大后缀，必须提高 pivot；从右找 successor 保证提高幅度最小，后缀反转保证整体刚好是下一个排列。

\textbf{边界实验。} 先用完全降序、重复数字、只有一个元素、交换后后缀必须最小化做反例检查。实验时把完全降序、重复数字、只有一个元素、交换后后缀必须最小化列成表，再打印关键下标和有效区间。若某个元素被写入后又被未来步骤破坏，说明区间不变量没有成立。

\textbf{复杂度变化。} 暴力阶乘级；原地算法三次线性扫描，时间 O(n)，空间 O(1)。

\textbf{迁移追问。} 如果题目改成“若要求上一个排列，比较方向和后缀处理对称变化”，先判断原状态是否仍然覆盖新答案集合，再决定是微调边界、改 value 语义，还是换算法。

\noindent\textbf{LeetCode 41 缺失的第一个正数。}

\textbf{题目说明。} LeetCode 41 缺失的第一个正数 的题面入口要先还原成具体任务：答案只可能落在一到 n 加一之间，数组本身可以充当值到位置的映射。

\textbf{样例入口。} 拿 [3, 4, -1, 1] 手算。答案只可能在 1 到 n+1；值 3 应放到下标 2，值 4 放到下标 3，-1 不管，1 放到下标 0。放完后第一个位置不等于 i+1 的地方就是缺失值。再试 [1, 1]，若不判断目标位置已有相同值会无限交换。

\textbf{暴力基线。} 暴力可以从 1 开始逐个查找是否出现在数组中，或者用哈希集合记录所有正数后再查。

\textbf{暴力过程。} 以 [3, 4, -1, 1] 为例，答案只可能在 1..5。把 3 放到下标 2，把 4 放到下标 3，把 1 放到下标 0，最后下标 1 缺 2。

\textbf{重复工作。} 题面模型：答案只可能落在一到 n 加一之间，数组本身可以充当值到位置的映射。暴力版的慢点要落到本题：浪费在额外集合或重复查找。长度为 n 的数组只关心 1..n 这些值，数组位置本身可以当哈希桶。后面的优化只消掉这类重复工作，不能改变答案集合。

\textbf{优化条件。} 本题的关键观察是：把合法正数 x 尽量交换到下标 x 减一，最后第一个位置不匹配就是答案。这句话必须能翻译成可维护状态；如果题目条件改变后观察不再成立，就不能继续套原来的写法。

\textbf{相邻状态。} 规律是数组本身当哈希表，只处理合法正数且避免重复交换。把这个特例继续拆开看：先标出已经确认的前缀或后缀，再标出未知区；能被丢掉的候选，必须是以后再也不会改变已确认区域语义的候选。这样才算从例子推出数组不变量，而不是只记一次操作。

\textbf{状态复用。} 扫描每个位置，把合法值 x 交换到 nums[x-1]。若目标位置已经是 x，停止交换，避免重复值死循环。手算时只改被新输入影响的那一小部分，而不是回头重算完整候选。

\textbf{指针和字段。} 状态字段要能说清：i 是当前整理位置，x=nums[i] 是试图归位的值，x-1 是它的目标下标。手算时按这个协议跟踪：[3, 4, -1, 1] 中 i=0，3 应去下标 2，换来 -1；i=1，4 应去下标 3，换来 1；1 再换到下标 0。

\textbf{代码落点。} 关键 while 条件是 1<=nums[i]<=n \&\& nums[nums[i]-1]!=nums[i]。第二遍找第一个 nums[i]!=i+1。关键代码行必须能逐句翻译成“旧状态怎样变成新状态”，否则就只是模板。

\textbf{正确性。} 每次有效交换都会至少让一个合法正数回到正确位置；最终若位置 i 不等于 i+1，则 i+1 没有被任何位置放回来，就是最小缺失正数。

\textbf{边界实验。} 先用非正数、大于 n 的数、重复值导致无限交换、数组已经完整包含一到 n做反例检查。实验时把非正数、大于 n 的数、重复值导致无限交换、数组已经完整包含一到 n列成表，再打印关键下标和有效区间。若某个元素被写入后又被未来步骤破坏，说明区间不变量没有成立。

\textbf{复杂度变化。} 哈希法 O(n) 空间；原地归位每个数交换有限次，总时间 O(n)，额外空间 O(1)。

\textbf{迁移追问。} 如果题目改成“若不能修改输入，只能使用哈希集合或额外布尔数组换取更简单边界”，先判断原状态是否仍然覆盖新答案集合，再决定是微调边界、改 value 语义，还是换算法。

\noindent\textbf{LeetCode 54 螺旋矩阵。}

\textbf{题目说明。} LeetCode 54 螺旋矩阵 的题面入口要先还原成具体任务：四个边界收缩模拟一圈一圈访问。

\textbf{样例入口。} 拿 3x3 矩阵手算，先走上边一整行，再走右边一整列，然后收缩 top 和 right；反向走下边和左边前要检查 top<=bottom、left<=right。

\textbf{暴力基线。} 暴力可以维护 visited 矩阵，按当前方向走，撞墙或访问过就转向，直到访问完所有格子。

\textbf{暴力过程。} 以 3x3 矩阵为例，先走上边一行，再走右边一列，再走下边和左边。最后只剩中心格时，如果不检查边界会重复访问。

\textbf{重复工作。} 题面模型：四个边界收缩模拟一圈一圈访问。暴力版的慢点要落到本题：浪费在 visited 矩阵。螺旋访问的未访问区域始终是一个收缩矩形，用四个边界就能表达。后面的优化只消掉这类重复工作，不能改变答案集合。

\textbf{优化条件。} 本题的关键观察是：每走完一条边就收缩对应边界，并在走反方向前检查是否仍合法。这句话必须能翻译成可维护状态；如果题目条件改变后观察不再成立，就不能继续套原来的写法。

\textbf{相邻状态。} 规律是四个边界每走完一条边就收缩一次，最后剩一行或一列时不能重复访问。把这个特例继续拆开看：先标出已经确认的前缀或后缀，再标出未知区；能被丢掉的候选，必须是以后再也不会改变已确认区域语义的候选。这样才算从例子推出数组不变量，而不是只记一次操作。

\textbf{状态复用。} 维护 top、bottom、left、right 四条边界。每走完一条边就收缩对应边界，反向走之前检查是否仍有行或列。手算时只改被新输入影响的那一小部分，而不是回头重算完整候选。

\textbf{指针和字段。} 状态字段要能说清：top/bottom 是当前剩余矩形的上下边，left/right 是左右边，res 保存访问顺序。手算时按这个协议跟踪：3x3 中 top 行访问完后 top++；right 列访问完后 right--；若 top<=bottom 再走底边，若 left<=right 再走左边。

\textbf{代码落点。} 关键是四段循环之间的边界检查，尤其是走下边和左边前要判断 top<=bottom、left<=right。关键代码行必须能逐句翻译成“旧状态怎样变成新状态”，否则就只是模板。

\textbf{正确性。} 每轮完整处理当前外圈，并把外圈从剩余矩形中删除；边界检查保证单行或单列时不会重复访问。

\textbf{边界实验。} 先用单行、单列、空矩阵、最后一圈只剩一个元素做反例检查。实验时把单行、单列、空矩阵、最后一圈只剩一个元素列成表，再打印关键下标和有效区间。若某个元素被写入后又被未来步骤破坏，说明区间不变量没有成立。

\textbf{复杂度变化。} visited 暴力 O(mn) 空间；边界法 O(mn) 时间，除输出外 O(1) 空间。

\textbf{迁移追问。} 如果题目改成“若改成生成矩阵，边界逻辑相同，只是读变成写”，先判断原状态是否仍然覆盖新答案集合，再决定是微调边界、改 value 语义，还是换算法。

\noindent\textbf{LeetCode 75 颜色分类。}

\textbf{题目说明。} LeetCode 75 颜色分类 的题面入口要先还原成具体任务：三指针维护零区、一号区、未知区和二区。

\textbf{样例入口。} 拿 [2, 0, 1] 手算，mid 指向 2 时和右边界交换，换回来的元素未知，所以 mid 不能前进；再看到 1 才前进。

\textbf{暴力基线。} 暴力可以计数 0、1、2 的个数后重写数组，或者排序数组。

\textbf{暴力过程。} 以 [2, 0, 1] 为例，计数法先数出每种颜色，再写 0、1、2。原地一趟要在扫描时把 0 放左边、2 放右边。

\textbf{重复工作。} 题面模型：三指针维护零区、一号区、未知区和二区。暴力版的慢点要落到本题：浪费在两趟计数或排序。颜色只有三类，可以在一趟扫描中维护四个区域。后面的优化只消掉这类重复工作，不能改变答案集合。

\textbf{优化条件。} 本题的关键观察是：遇到二时只移动右边界，不移动扫描指针，因为换回来的元素未知。这句话必须能翻译成可维护状态；如果题目条件改变后观察不再成立，就不能继续套原来的写法。

\textbf{相邻状态。} 规律是 [0, low) 是 0， [low, mid) 是 1， (high, n) 是 2，未知区只在 mid 到 high 之间收缩。把这个特例继续拆开看：先标出已经确认的前缀或后缀，再标出未知区；能被丢掉的候选，必须是以后再也不会改变已确认区域语义的候选。这样才算从例子推出数组不变量，而不是只记一次操作。

\textbf{状态复用。} 维护 low、mid、high：[0, low) 是 0，[low, mid) 是 1，[mid, high] 是未知，(high, n) 是 2。手算时只改被新输入影响的那一小部分，而不是回头重算完整候选。

\textbf{指针和字段。} 状态字段要能说清：mid 是当前检查位置，low 是下一个 0 的位置，high 是下一个 2 的位置。手算时按这个协议跟踪：[2, 0, 1] 中 mid 指向 2，和 high 交换后 high--，mid 不动，因为换回来的 1 还没检查；再处理 1，mid++。

\textbf{代码落点。} 关键是遇到 0 时 swap(nums[low++], nums[mid++])；遇到 2 时 swap(nums[mid], nums[high--])，mid 不能前进。关键代码行必须能逐句翻译成“旧状态怎样变成新状态”，否则就只是模板。

\textbf{正确性。} 每次操作都把一个元素放入 0 区或 2 区，或确认它属于 1 区；未知区严格缩小，不变量保持。

\textbf{边界实验。} 先用全零、全二、已经有序、逆序、交换后漏检查做反例检查。实验时把全零、全二、已经有序、逆序、交换后漏检查列成表，再打印关键下标和有效区间。若某个元素被写入后又被未来步骤破坏，说明区间不变量没有成立。

\textbf{复杂度变化。} 排序 O(n log n)，计数 O(n) 两趟；荷兰国旗算法 O(n) 时间、O(1) 空间。

\textbf{迁移追问。} 如果题目改成“这是荷兰国旗问题，能推广到三类分区的原地维护”，先判断原状态是否仍然覆盖新答案集合，再决定是微调边界、改 value 语义，还是换算法。

\noindent\textbf{LeetCode 88 合并两个有序数组。}

\textbf{题目说明。} LeetCode 88 合并两个有序数组 的题面入口要先还原成具体任务：第一个数组尾部有空位，适合从后往前写入最终最大值。

\textbf{样例入口。} 拿 nums1=[1, 3, 0]、nums2=[2] 手算，如果从前往后写，2 会覆盖 nums1 中尚未比较的 3；从尾部写，先比较 3 和 2，把 3 放到最后，再把 2 放到中间。

\textbf{暴力基线。} 暴力可以把 nums2 追加到 nums1 有效元素后再排序，或者新建数组从前往后归并。

\textbf{暴力过程。} 以 nums1=[1, 3, 0]、nums2=[2] 为例，从前往后写会把 nums1 中尚未读取的 3 覆盖掉。

\textbf{重复工作。} 题面模型：第一个数组尾部有空位，适合从后往前写入最终最大值。暴力版的慢点要落到本题：浪费在排序或额外数组。nums1 尾部本来有空位，应该从后往前写入当前最大值，避免覆盖未读元素。后面的优化只消掉这类重复工作，不能改变答案集合。

\textbf{优化条件。} 本题的关键观察是：逆向填充避免覆盖尚未读取的 nums1 元素。这句话必须能翻译成可维护状态；如果题目条件改变后观察不再成立，就不能继续套原来的写法。

\textbf{相邻状态。} 规律是利用尾部空位逆向填充，避免覆盖未读元素。把这个特例继续拆开看：先标出已经确认的前缀或后缀，再标出未知区；能被丢掉的候选，必须是以后再也不会改变已确认区域语义的候选。这样才算从例子推出数组不变量，而不是只记一次操作。

\textbf{状态复用。} i=m-1 指向 nums1 有效尾，j=n-1 指向 nums2 尾，write=m+n-1 指向最终写入位置。手算时只改被新输入影响的那一小部分，而不是回头重算完整候选。

\textbf{指针和字段。} 状态字段要能说清：i/j 是两个有序数组尚未合并部分的尾部，write 是结果数组尚未填充的最后位置。手算时按这个协议跟踪：[1, 3, 0] 和 [2] 中，比较 3 和 2，把 3 写到 write=2；再把 2 写到 write=1，1 已在原位。

\textbf{代码落点。} 关键代码是从尾部 while j>=0 合并；若 nums1 剩余不用搬，若 nums2 剩余必须继续写入。关键代码行必须能逐句翻译成“旧状态怎样变成新状态”，否则就只是模板。

\textbf{正确性。} 最终数组最大元素一定来自两个当前尾部之一，把它写到最后不会影响剩余元素的相对有序性。

\textbf{边界实验。} 先用第二个数组为空、第一个有效元素为空、重复值、尾部剩余做反例检查。实验时把第二个数组为空、第一个有效元素为空、重复值、尾部剩余列成表，再打印关键下标和有效区间。若某个元素被写入后又被未来步骤破坏，说明区间不变量没有成立。

\textbf{复杂度变化。} 排序 O((m+n)log(m+n))；逆向归并 O(m+n) 时间、O(1) 额外空间。

\textbf{迁移追问。} 如果题目改成“若没有额外空位，只能新开数组或使用更复杂原地归并”，先判断原状态是否仍然覆盖新答案集合，再决定是微调边界、改 value 语义，还是换算法。

\noindent\textbf{LeetCode 189 轮转数组。}

\textbf{题目说明。} LeetCode 189 轮转数组给定数组 nums 和整数 k，要求把数组原地向右轮转 k 个位置。

\textbf{样例入口。} 样例 nums=[1, 2, 3, 4, 5, 6, 7]，k=3，轮转后为 [5, 6, 7, 1, 2, 3, 4]。

\textbf{暴力基线。} 暴力可以新建 answer，让 answer[(i+k)\%n]=nums[i]，最后复制回 nums。

\textbf{暴力过程。} 以 [1, 2, 3, 4, 5, 6, 7]、k=3 为例，1 去下标 3，5 去下标 0，结果是 [5, 6, 7, 1, 2, 3, 4]。

\textbf{重复工作。} 题面模型：给定数组和整数 k，将数组中的元素整体向右轮转 k 个位置。暴力版的慢点要落到本题：浪费在额外数组。右移 k 位等价于把后 k 个元素整体搬到前面，原地可以用反转实现两段换序。后面的优化只消掉这类重复工作，不能改变答案集合。

\textbf{优化条件。} 本题的关键观察是：右移 k 位等价于把数组切成两段后换序，三次反转能原地完成：整体、前 k、后 n 减 k。这句话必须能翻译成可维护状态；如果题目条件改变后观察不再成立，就不能继续套原来的写法。

\textbf{相邻状态。} 规律是右移 k 等价于两段换序，k 要先对 n 取模。把这个特例继续拆开看：先标出已经确认的前缀或后缀，再标出未知区；能被丢掉的候选，必须是以后再也不会改变已确认区域语义的候选。这样才算从例子推出数组不变量，而不是只记一次操作。

\textbf{状态复用。} 先 k\%=n；整体反转让后段到前面，再反转前 k 段和后 n-k 段，恢复每段内部顺序。手算时只改被新输入影响的那一小部分，而不是回头重算完整候选。

\textbf{指针和字段。} 状态字段要能说清：k 是有效轮转距离，三个区间分别是 whole、[0, k)、[k, n)。手算时按这个协议跟踪：[1, 2, 3, 4, 5, 6, 7] 整体反转为 [7, 6, 5, 4, 3, 2, 1]；前 3 个反转为 [5, 6, 7, 4, 3, 2, 1]；后 4 个反转为 [5, 6, 7, 1, 2, 3, 4]。

\textbf{代码落点。} 关键是空数组先返回，否则 k 取模会除零。随后按“整体反转、修正前段、修正后段”三步写，不把长函数调用塞成一整行。关键代码行必须能逐句翻译成“旧状态怎样变成新状态”，否则就只是模板。

\textbf{正确性。} 整体反转完成两段位置交换但段内顺序反了；再分别反转两段，正好得到右移后的段顺序。

\textbf{边界实验。} 先用k 为零、k 大于 n、空数组、单元素做反例检查。实验时把k 为零、k 大于 n、空数组、单元素列成表，再打印关键下标和有效区间。若某个元素被写入后又被未来步骤破坏，说明区间不变量没有成立。

\textbf{复杂度变化。} 辅助数组 O(n) 空间；三次反转 O(n) 时间、O(1) 额外空间。

\textbf{迁移追问。} 如果题目改成“环状替换也可做，但要处理最大公约数个循环”，先判断原状态是否仍然覆盖新答案集合，再决定是微调边界、改 value 语义，还是换算法。

\noindent\textbf{LeetCode 204 计数质数。}

\textbf{题目说明。} LeetCode 204 计数质数 的题面入口要先还原成具体任务：逐个数试除会重复判断很多合数的因子。

\textbf{样例入口。} 拿 n=10 手算，2 是质数，标记 4、6、8；3 是质数，标记 9；4 已被标记跳过，剩下质数是 2、3、5、7。

\textbf{暴力基线。} 暴力逐个判断 2..n-1 是否为质数，每个数再试除 2..sqrt(x)。

\textbf{暴力过程。} 以 n=10 为例，暴力会分别判断 4、6、8 是否能被 2 整除，也会在多个合数上重复试同一个因子。

\textbf{重复工作。} 题面模型：逐个数试除会重复判断很多合数的因子。暴力版的慢点要落到本题：浪费在反复证明合数是合数。一个数一旦被某个质数标记为倍数，后续不必再试除判断。后面的优化只消掉这类重复工作，不能改变答案集合。

\textbf{优化条件。} 本题的关键观察是：埃氏筛从每个质数的平方开始标记它的倍数，未被标记的数就是质数。这句话必须能翻译成可维护状态；如果题目条件改变后观察不再成立，就不能继续套原来的写法。

\textbf{相邻状态。} 规律是埃氏筛只从质数的平方开始标记倍数，避免小倍数被重复处理。把这个特例继续拆开看：先标出已经确认的前缀或后缀，再标出未知区；能被丢掉的候选，必须是以后再也不会改变已确认区域语义的候选。这样才算从例子推出数组不变量，而不是只记一次操作。

\textbf{状态复用。} is\_prime 或 is\_composite 数组记录标记状态。从质数 i 的 i*i 开始标记 i 的倍数。手算时只改被新输入影响的那一小部分，而不是回头重算完整候选。

\textbf{指针和字段。} 状态字段要能说清：i 是当前候选质数，multiple 是被标记的合数，count 是未被标记的质数数量。手算时按这个协议跟踪：n=10 时，i=2 标记 4、6、8；i=3 从 9 开始标记；4 已被标记跳过，质数为 2、3、5、7。

\textbf{代码落点。} 关键是 for (long long j=1LL*i*i; j<n; j+=i) 标记，边界是不包含 n。关键代码行必须能逐句翻译成“旧状态怎样变成新状态”，否则就只是模板。

\textbf{正确性。} 小于 i*i 的 i 的倍数已经有更小因子标记过；未被任何小因子标记的数只能是质数。

\textbf{边界实验。} 先用n 小于二、从 i*i 开始防溢出、重复标记、边界是不包含 n做反例检查。实验时把n 小于二、从 i*i 开始防溢出、重复标记、边界是不包含 n列成表，再打印关键下标和有效区间。若某个元素被写入后又被未来步骤破坏，说明区间不变量没有成立。

\textbf{复杂度变化。} 试除法约 O(n sqrt n)，埃氏筛 O(n log log n) 时间、O(n) 空间。

\textbf{迁移追问。} 如果题目改成“若要多次查询前缀质数个数，可以在筛后再做前缀计数”，先判断原状态是否仍然覆盖新答案集合，再决定是微调边界、改 value 语义，还是换算法。

\noindent\textbf{LeetCode 238 除自身以外数组的乘积。}

\textbf{题目说明。} LeetCode 238 除自身以外数组的乘积给定整数数组 nums，返回 answer，其中 answer[i] 等于 nums 中除 nums[i] 以外所有元素的乘积，且不能使用除法。

\textbf{样例入口。} 样例 nums=[1, 2, 3, 4]，输出 [24, 12, 8, 6]；每个位置都排除自身后乘其余元素。

\textbf{暴力基线。} 暴力对每个位置 i 扫描整个数组，跳过 i，把其他元素全部相乘得到 answer[i]。

\textbf{暴力过程。} 以 [1, 2, 3, 4] 为例，answer[2] 要算 1*2*4；answer[3] 又重算 1*2*3，左侧乘积被重复计算。

\textbf{重复工作。} 题面模型：给定整数数组，返回 answer，其中 answer[i] 等于 nums[i] 之外所有元素的乘积。暴力版的慢点要落到本题：浪费在每个 i 都重算左侧和右侧乘积。位置 i 的答案天然等于 i 左侧乘积乘 i 右侧乘积。后面的优化只消掉这类重复工作，不能改变答案集合。

\textbf{优化条件。} 本题的关键观察是：不能用除法时，答案由左侧乘积和右侧乘积组成；先写入左侧前缀积，再从右向左乘右侧后缀积。这句话必须能翻译成可维护状态；如果题目条件改变后观察不再成立，就不能继续套原来的写法。

\textbf{相邻状态。} 规律是答案等于左侧乘积乘右侧乘积，不能使用除法时就用两次扫描合并贡献。把这个特例继续拆开看：先标出已经确认的前缀或后缀，再标出未知区；能被丢掉的候选，必须是以后再也不会改变已确认区域语义的候选。这样才算从例子推出数组不变量，而不是只记一次操作。

\textbf{状态复用。} 第一遍从左到右，把 answer[i] 写成 i 左侧所有数的乘积；第二遍从右到左，用 right\_product 乘回右侧贡献。手算时只改被新输入影响的那一小部分，而不是回头重算完整候选。

\textbf{指针和字段。} 状态字段要能说清：left\_product 隐含在 answer[i]，right\_product 是当前 i 右侧乘积，answer[i] 最终是左右贡献相乘。手算时按这个协议跟踪：[1, 2, 3, 4] 第一遍 answer=[1, 1, 2, 6]；第二遍 right\_product 依次为 1、4、12、24，最终 [24, 12, 8, 6]。

\textbf{代码落点。} 关键是第二遍先 answer[i]*=right\_product，再 right\_product*=nums[i]，这样不会把 nums[i] 自己乘进去。关键代码行必须能逐句翻译成“旧状态怎样变成新状态”，否则就只是模板。

\textbf{正确性。} 第一遍保证 answer[i] 只含左侧元素，第二遍乘入的 right\_product 只含右侧元素，二者乘积正好排除自身。

\textbf{边界实验。} 先用一个零、多个零、负数、乘积溢出做反例检查。实验时把一个零、多个零、负数、乘积溢出列成表，再打印关键下标和有效区间。若某个元素被写入后又被未来步骤破坏，说明区间不变量没有成立。

\textbf{复杂度变化。} 暴力 O(\ensuremath{n^{2}})，前后缀合并 O(n) 时间；若输出数组不计额外空间，则额外空间 O(1)。

\textbf{迁移追问。} 如果题目改成“这个题展示输出数组可作为状态存储的一部分”，先判断原状态是否仍然覆盖新答案集合，再决定是微调边界、改 value 语义，还是换算法。

\noindent\textbf{LeetCode 283 移动零。}

\textbf{题目说明。} LeetCode 283 移动零给定数组 nums，要求原地把所有 0 移到末尾，同时保持非零元素的相对顺序。

\textbf{样例入口。} 样例 nums=[0, 1, 0, 3, 12]，处理后为 [1, 3, 12, 0, 0]；非零元素 1、3、12 的相对顺序不变。

\textbf{暴力基线。} 暴力可以新建数组，先放所有非零元素，再补同样数量的零，最后复制回原数组。

\textbf{暴力过程。} 以 [0, 1, 0, 3, 12] 为例，辅助数组先收集 [1, 3, 12]，再补 [0, 0]。

\textbf{重复工作。} 题面模型：给定数组，原地把所有 0 移到末尾，同时保持非零元素的相对顺序。暴力版的慢点要落到本题：浪费在辅助数组。非零元素的相对顺序由读取顺序天然保证，只需要一个原地写入位置。后面的优化只消掉这类重复工作，不能改变答案集合。

\textbf{优化条件。} 本题的关键观察是：稳定保留非零元素顺序，把零集中到尾部；慢指针写入下一个非零位置，最后补零或用交换。这句话必须能翻译成可维护状态；如果题目条件改变后观察不再成立，就不能继续套原来的写法。

\textbf{相邻状态。} 规律是先稳定压缩非零前缀，再填零；不能在扫描时盲目交换导致顺序变化。把这个特例继续拆开看：先标出已经确认的前缀或后缀，再标出未知区；能被丢掉的候选，必须是以后再也不会改变已确认区域语义的候选。这样才算从例子推出数组不变量，而不是只记一次操作。

\textbf{状态复用。} write 指向下一个非零元素应写入的位置。第一遍把非零元素稳定写到前面，第二遍把 write 之后填 0。手算时只改被新输入影响的那一小部分，而不是回头重算完整候选。

\textbf{指针和字段。} 状态字段要能说清：read 是扫描位置，write 左侧是不含 0 的稳定前缀，write 到末尾最后会被置零。手算时按这个协议跟踪：读到 1 写到 nums[0]，读到 3 写到 nums[1]，读到 12 写到 nums[2]，最后 nums[3]、nums[4] 置 0。

\textbf{代码落点。} 关键代码是 if (nums[read]!=0) nums[write++]=nums[read]; 后面 while(write<n) nums[write++]=0。关键代码行必须能逐句翻译成“旧状态怎样变成新状态”，否则就只是模板。

\textbf{正确性。} 所有非零元素按原读取顺序写入前缀，因此稳定性保持；剩余位置数量等于原零的数量，补零后满足题意。

\textbf{边界实验。} 先用全零、无零、零在开头结尾、稳定顺序要求做反例检查。实验时把全零、无零、零在开头结尾、稳定顺序要求列成表，再打印关键下标和有效区间。若某个元素被写入后又被未来步骤破坏，说明区间不变量没有成立。

\textbf{复杂度变化。} 辅助数组 O(n) 空间；原地稳定写 O(n) 时间、O(1) 额外空间。

\textbf{迁移追问。} 如果题目改成“若目标是移动指定值，模型与移除元素相同”，先判断原状态是否仍然覆盖新答案集合，再决定是微调边界、改 value 语义，还是换算法。

\noindent\textbf{LeetCode 304 二维区域和检索 - 矩阵不可变。}

\textbf{题目说明。} LeetCode 304 二维区域和检索 - 矩阵不可变给定一个固定矩阵，先构造 NumMatrix，再多次调用 sumRegion(row1, col1, row2, col2) 查询闭矩形区域和。

\textbf{样例入口。} 题面样例的三次查询分别返回 8、11、12；正文只保留查询结果，完整矩阵见原题样例。

\textbf{暴力基线。} 暴力每次查询都双层循环扫描 row1..row2、col1..col2，把矩形内元素相加。

\textbf{暴力过程。} 若多次查询同一个固定矩阵，sumRegion(2, 1, 4, 3) 会扫描一块矩形，下一次相邻矩形又会重复扫描大量格子。

\textbf{重复工作。} 题面模型：固定矩阵上需要多次查询闭矩形区域和。暴力版的慢点要落到本题：浪费在矩阵不变却每次重新读格子。任意查询矩形都可以由几个左上前缀矩形组合出来。后面的优化只消掉这类重复工作，不能改变答案集合。

\textbf{优化条件。} 本题的关键观察是：预处理二维前缀和后，每次查询用四个角做容斥。这句话必须能翻译成可维护状态；如果题目条件改变后观察不再成立，就不能继续套原来的写法。

\textbf{相邻状态。} 规律是查询公式对应“右下、上方、左侧、左上重叠”四个量；多开一行一列哨兵可以统一第一行和第一列。把这个特例继续拆开看：先标出已经确认的前缀或后缀，再标出未知区；能被丢掉的候选，必须是以后再也不会改变已确认区域语义的候选。这样才算从例子推出数组不变量，而不是只记一次操作。

\textbf{状态复用。} 预处理 prefix[r+1][c+1] 表示原矩阵 [0..r][0..c] 的和。查询时用右下大矩形减上方、减左方、加左上重叠。手算时只改被新输入影响的那一小部分，而不是回头重算完整候选。

\textbf{指针和字段。} 状态字段要能说清：prefix 的多一行多一列是哨兵，row1/col1/row2/col2 是原矩阵闭区间坐标，查询公式使用加一后的前缀坐标。手算时按这个协议跟踪：sumRegion(2, 1, 4, 3) 等于 P[5][4]-P[2][4]-P[5][1]+P[2][1]。

\textbf{代码落点。} 构建时分四项：当前格、上方前缀、左方前缀，再减左上重叠。代码里分别对应 matrix 当前值、上方 prefix、左方 prefix 和左上 prefix。关键代码行必须能逐句翻译成“旧状态怎样变成新状态”，否则就只是模板。

\textbf{正确性。} 右下大矩形包含目标、上方和左方；上方左方相减时左上角被减了两次，所以要加回一次。

\textbf{边界实验。} 先用第一行第一列、单行单列、空矩阵、多次重复查询、坐标加一做反例检查。实验时把第一行第一列、单行单列、空矩阵、多次重复查询、坐标加一列成表，再打印关键下标和有效区间。若某个元素被写入后又被未来步骤破坏，说明区间不变量没有成立。

\textbf{复杂度变化。} 单次暴力 O(矩形面积)，预处理 O(mn)，每次查询 O(1)，空间 O(mn)。

\textbf{迁移追问。} 如果题目改成“若矩阵会更新，普通二维前缀和失效，需要二维树状数组或线段树”，先判断原状态是否仍然覆盖新答案集合，再决定是微调边界、改 value 语义，还是换算法。

\noindent\textbf{LeetCode 307 区域和检索 - 数组可修改。}

\textbf{题目说明。} LeetCode 307 区域和检索 - 数组可修改 的题面入口要先还原成具体任务：单点更新会破坏普通前缀和后面所有位置。

\textbf{样例入口。} 拿 nums=[1, 3, 5] 手算，sumRange(0, 2)=9；update(1, 2) 后数组变成 [1, 2, 5]，后续区间和都受影响。普通前缀和要改很多位置，树状数组只把差值 -1 加到若干 lowbit 管辖节点。

\textbf{暴力基线。} 普通前缀和能 O(1) 查询，但 update(index, val) 后 index 之后所有前缀都要修改。

\textbf{暴力过程。} 以 nums=[1, 3, 5] 为例，update(1, 2) 让数组变成 [1, 2, 5]。若维护 prefix，prefix[2]、prefix[3] 都要改。

\textbf{重复工作。} 题面模型：单点更新会破坏普通前缀和后面所有位置。暴力版的慢点要落到本题：浪费在一次单点更新影响大量前缀。树状数组把前缀和拆成若干 lowbit 管辖段，更新只触碰包含该点的少数段。后面的优化只消掉这类重复工作，不能改变答案集合。

\textbf{优化条件。} 本题的关键观察是：树状数组把前缀和拆成若干 lowbit 管辖段，更新和查询都只访问对数个节点。这句话必须能翻译成可维护状态；如果题目条件改变后观察不再成立，就不能继续套原来的写法。

\textbf{相邻状态。} 规律是树状数组用对数个节点维护前缀和，更新和查询都沿 lowbit 跳转。把这个特例继续拆开看：先标出已经确认的前缀或后缀，再标出未知区；能被丢掉的候选，必须是以后再也不会改变已确认区域语义的候选。这样才算从例子推出数组不变量，而不是只记一次操作。

\textbf{状态复用。} BIT 使用 1 基下标。update 计算 delta，然后 for i+=lowbit(i) 更新管辖段；query 前缀和时 for i-=lowbit(i) 累加。手算时只改被新输入影响的那一小部分，而不是回头重算完整候选。

\textbf{指针和字段。} 状态字段要能说清：tree[i] 保存一段长度 lowbit(i) 的区间和，nums 保存当前真实值以计算 delta。手算时按这个协议跟踪：update(1, 2) 的 1 基位置是 2，delta=-1，更新 tree[2]、tree[4] 等管辖节点。

\textbf{代码落点。} 关键代码是 index += index \& -index 和 index -= index \& -index；sumRange(l, r)=prefix(r+1)-prefix(l)。关键代码行必须能逐句翻译成“旧状态怎样变成新状态”，否则就只是模板。

\textbf{正确性。} 任意前缀可以被若干 lowbit 段不重叠拆分；单点更新只需更新所有包含该点的拆分段。

\textbf{边界实验。} 先用下标一基转换、负数更新、重复更新、查询单点、差值更新做反例检查。实验时把下标一基转换、负数更新、重复更新、查询单点、差值更新列成表，再打印关键下标和有效区间。若某个元素被写入后又被未来步骤破坏，说明区间不变量没有成立。

\textbf{复杂度变化。} 普通前缀 update O(n)，树状数组 update/query 均 O(log n)，空间 O(n)。

\textbf{迁移追问。} 如果题目改成“若需要区间更新或更复杂信息，线段树更通用”，先判断原状态是否仍然覆盖新答案集合，再决定是微调边界、改 value 语义，还是换算法。

\noindent\textbf{LeetCode 344 反转字符串。}

\textbf{题目说明。} LeetCode 344 反转字符串 的题面入口要先还原成具体任务：反转只需要把对称位置交换，没必要新建完整副本。

\textbf{样例入口。} 拿 s=[h, e, l, l, o] 手算，交换首尾得到 o, e, l, l, h，再交换 e 和 l，最终 o, l, l, e, h。

\textbf{暴力基线。} 暴力可以新建一个反向数组，把原数组从后往前复制过去，再写回。

\textbf{暴力过程。} 以 [h, e, l, l, o] 为例，辅助数组得到 [o, l, l, e, h]。原地做法只需要交换首尾、次首尾。

\textbf{重复工作。} 题面模型：反转只需要把对称位置交换，没必要新建完整副本。暴力版的慢点要落到本题：浪费在完整副本。反转的每个位置只和它的对称位置交换一次。后面的优化只消掉这类重复工作，不能改变答案集合。

\textbf{优化条件。} 本题的关键观察是：左右指针从两端向中间移动，每轮交换一对字符。这句话必须能翻译成可维护状态；如果题目条件改变后观察不再成立，就不能继续套原来的写法。

\textbf{相邻状态。} 规律是左右指针每轮交换对称字符并向中间收缩，原地反转不需要额外数组。把这个特例继续拆开看：先标出已经确认的前缀或后缀，再标出未知区；能被丢掉的候选，必须是以后再也不会改变已确认区域语义的候选。这样才算从例子推出数组不变量，而不是只记一次操作。

\textbf{状态复用。} left 从 0 开始，right 从 n-1 开始；left<right 时交换 s[left] 和 s[right]，然后向中间收缩。手算时只改被新输入影响的那一小部分，而不是回头重算完整候选。

\textbf{指针和字段。} 状态字段要能说清：left/right 是尚未交换的对称位置，二者之间是还未完全处理的区域。手算时按这个协议跟踪：先交换 h 和 o，得到 [o, e, l, l, h]；再交换 e 和 l，得到 [o, l, l, e, h]。

\textbf{代码落点。} 关键代码是 while (left < right) swap(s[left++], s[right--])。关键代码行必须能逐句翻译成“旧状态怎样变成新状态”，否则就只是模板。

\textbf{正确性。} 每次交换把一对对称位置放到最终位置，已处理两端不会再被修改，未知区逐步缩小。

\textbf{边界实验。} 先用空数组、单字符、偶数长度、奇数长度、原地修改做反例检查。实验时把空数组、单字符、偶数长度、奇数长度、原地修改列成表，再打印关键下标和有效区间。若某个元素被写入后又被未来步骤破坏，说明区间不变量没有成立。

\textbf{复杂度变化。} 辅助数组 O(n) 空间；双指针 O(n) 时间、O(1) 额外空间。

\textbf{迁移追问。} 如果题目改成“若反转单词顺序，要先定义分隔符和单词边界”，先判断原状态是否仍然覆盖新答案集合，再决定是微调边界、改 value 语义，还是换算法。

\noindent\textbf{LeetCode 459 重复的子字符串。}

\textbf{题目说明。} LeetCode 459 重复的子字符串 的题面入口要先还原成具体任务：判断字符串能否由某个较短周期重复组成，本质是找合法周期长度。

\textbf{样例入口。} 拿 s=abab 手算，最短周期 ab 重复两次；拿 aba 手算，虽然有前后缀 a，但剩余周期长度 2 不能整除 3。

\textbf{暴力基线。} 暴力枚举可能周期 len，若 n 能被 len 整除，就逐段检查 s 是否由 s[0..len-1] 重复组成。

\textbf{暴力过程。} 以 s=abab 为例，len=1 失败，len=2 的周期 ab 重复两次成功；s=aba 虽有前后缀 a，但周期长度 2 不能整除 3。

\textbf{重复工作。} 题面模型：判断字符串能否由某个较短周期重复组成，本质是找合法周期长度。暴力版的慢点要落到本题：浪费在每个周期都重新比较。KMP 前缀函数已经记录了最长相等前后缀，从而推出最小可能周期。后面的优化只消掉这类重复工作，不能改变答案集合。

\textbf{优化条件。} 本题的关键观察是：可用 KMP 前缀函数得到最长相等前后缀，若剩余周期长度能整除总长则成立。这句话必须能翻译成可维护状态；如果题目条件改变后观察不再成立，就不能继续套原来的写法。

\textbf{相邻状态。} 规律是 KMP 前缀函数给出最长相等前后缀，周期长度为 n-lps，能整除才说明由子串重复组成。把这个特例继续拆开看：先标出已经确认的前缀或后缀，再标出未知区；能被丢掉的候选，必须是以后再也不会改变已确认区域语义的候选。这样才算从例子推出数组不变量，而不是只记一次操作。

\textbf{状态复用。} 计算 lps[n-1]，令 period=n-lps[n-1]。若 lps>0 且 n\%period==0，则可以由长度 period 的子串重复。手算时只改被新输入影响的那一小部分，而不是回头重算完整候选。

\textbf{指针和字段。} 状态字段要能说清：lps[i] 是 s[0..i] 的最长相等真前后缀长度，period 是候选重复单元长度。手算时按这个协议跟踪：abab 的 lps 最后为 2，period=2，4\%2==0，所以成立；aba 的 lps 最后为 1，period=2，3\%2!=0，所以不成立。

\textbf{代码落点。} 关键是构建前缀函数时失配用 j=lps[j-1] 回退；最终还要检查整除条件，不能只看 lps>0。关键代码行必须能逐句翻译成“旧状态怎样变成新状态”，否则就只是模板。

\textbf{正确性。} 若字符串由周期重复组成，则存在长度 n-period 的相等前后缀；反过来，最长前后缀推出的 period 能整除 n 时，整串可按该周期铺满。

\textbf{边界实验。} 先用长度一、全相同字符、无重复、部分前后缀相同但不能整除做反例检查。实验时把长度一、全相同字符、无重复、部分前后缀相同但不能整除列成表，再打印关键下标和有效区间。若某个元素被写入后又被未来步骤破坏，说明区间不变量没有成立。

\textbf{复杂度变化。} 枚举周期最坏 O(\ensuremath{n^{2}})，KMP 前缀函数 O(n) 时间、O(n) 空间。

\textbf{迁移追问。} 如果题目改成“也可用 s+s 去掉首尾后查找原串，但要能解释为什么排除了自身匹配”，先判断原状态是否仍然覆盖新答案集合，再决定是微调边界、改 value 语义，还是换算法。

\noindent\textbf{LeetCode 493 翻转对。}

\textbf{题目说明。} LeetCode 493 翻转对 的题面入口要先还原成具体任务：每个 i 需要统计右侧满足 nums[i] 大于 2*nums[j] 的元素。

\textbf{样例入口。} 拿 nums=[1, 3, 2, 3, 1] 手算，翻转对有 (3, 1)、(3, 1) 两个。暴力每个 i 扫右侧是平方级；归并排序中左右半区已有序，可以用指针统计左侧每个数能压过右侧多少个 2*nums[j]。

\textbf{暴力基线。} 暴力枚举所有 i<j，检查 nums[i] > 2*nums[j]，计数满足条件的对。

\textbf{暴力过程。} 以 [1, 3, 2, 3, 1] 为例，暴力会让每个左侧元素扫描右侧所有元素，找到两个翻转对。

\textbf{重复工作。} 题面模型：每个 i 需要统计右侧满足 nums[i] 大于 2*nums[j] 的元素。暴力版的慢点要落到本题：浪费在右侧反复扫描。归并排序中左右半区已经有序，可以用指针一次统计每个左元素能匹配多少右元素。后面的优化只消掉这类重复工作，不能改变答案集合。

\textbf{优化条件。} 本题的关键观察是：归并排序在合并前用两个有序半区双指针计数，再正常归并保持有序。这句话必须能翻译成可维护状态；如果题目条件改变后观察不再成立，就不能继续套原来的写法。

\textbf{相邻状态。} 规律是计数发生在归并前，随后归并保持区间有序。把这个特例继续拆开看：先标出已经确认的前缀或后缀，再标出未知区；能被丢掉的候选，必须是以后再也不会改变已确认区域语义的候选。这样才算从例子推出数组不变量，而不是只记一次操作。

\textbf{状态复用。} 分治排序。合并前，left 半区和 right 半区都有序；对每个 i，推进 j 直到 nums[i] <= 2*nums[j]，新增 j-mid-1 个对。手算时只改被新输入影响的那一小部分，而不是回头重算完整候选。

\textbf{指针和字段。} 状态字段要能说清：i 是左半区位置，j 是右半区第一个不满足的位置，merge 负责保持当前区间有序。手算时按这个协议跟踪：[1, 3, 2, 3, 1] 分治后，在合并有序半区时统计 3 与右侧 1 的关系，严格大于 2*1 成立。

\textbf{代码落点。} 关键是比较 1LL*nums[i] > 2LL*nums[j]，避免乘以 2 溢出；计数完成后仍要正常归并排序。关键代码行必须能逐句翻译成“旧状态怎样变成新状态”，否则就只是模板。

\textbf{正确性。} 右半区有序时，对固定 i，满足条件的右侧下标形成前缀；j 单调前进不会漏计也不会重复计。

\textbf{边界实验。} 先用负数、乘以二溢出、重复值、严格大于、答案数量很大做反例检查。实验时把负数、乘以二溢出、重复值、严格大于、答案数量很大列成表，再打印关键下标和有效区间。若某个元素被写入后又被未来步骤破坏，说明区间不变量没有成立。

\textbf{复杂度变化。} 暴力 O(\ensuremath{n^{2}})，归并计数 O(n log n)，辅助空间 O(n)。

\textbf{迁移追问。} 如果题目改成“也可用离散化加树状数组，关键仍是动态排名计数”，先判断原状态是否仍然覆盖新答案集合，再决定是微调边界、改 value 语义，还是换算法。

\noindent\textbf{LeetCode 581 最短无序连续子数组。}

\textbf{题目说明。} LeetCode 581 最短无序连续子数组 的题面入口要先还原成具体任务：要找破坏全局有序性的最左和最右位置。

\textbf{样例入口。} 拿 [2, 6, 4, 8, 10, 9, 15] 手算，从左扫描时 4 小于此前最大 6，说明右边界至少到 4；9 小于此前最大 10，右边界更新到 9。

\textbf{暴力基线。} 暴力可以复制数组排序，再比较原数组和排序数组的首尾不同位置，得到需要排序的最短区间。

\textbf{暴力过程。} 以 [2, 6, 4, 8, 10, 9, 15] 为例，排序后是 [2, 4, 6, 8, 9, 10, 15]，首个不同在 1，最后不同在 5。

\textbf{重复工作。} 题面模型：要找破坏全局有序性的最左和最右位置。暴力版的慢点要落到本题：浪费在完整排序。我们只需要找哪些位置破坏全局有序性，而不是得到完整有序数组。后面的优化只消掉这类重复工作，不能改变答案集合。

\textbf{优化条件。} 本题的关键观察是：从左维护历史最大值确定右边界，从右维护历史最小值确定左边界。这句话必须能翻译成可维护状态；如果题目条件改变后观察不再成立，就不能继续套原来的写法。

\textbf{相邻状态。} 规律是左右扫描分别找破坏有序的位置，最后得到最短需要排序的连续段。把这个特例继续拆开看：先标出已经确认的前缀或后缀，再标出未知区；能被丢掉的候选，必须是以后再也不会改变已确认区域语义的候选。这样才算从例子推出数组不变量，而不是只记一次操作。

\textbf{状态复用。} 从左到右维护历史最大值 max\_seen，若 nums[i]<max\_seen，说明 i 必须纳入右边界；从右到左维护历史最小值 min\_seen，若 nums[i]>min\_seen，说明 i 必须纳入左边界。手算时只改被新输入影响的那一小部分，而不是回头重算完整候选。

\textbf{指针和字段。} 状态字段要能说清：right 是最后一个小于历史最大值的位置，left 是最后一个大于右侧最小值的位置。手算时按这个协议跟踪：样例中从左扫到 4 时，小于历史最大 6，right=2；扫到 9 时，小于历史最大 10，right=5。反向扫可得到 left=1。

\textbf{代码落点。} 关键是两次扫描分别更新 max\_seen/min\_seen 和边界；若 right 没被更新，数组已经有序，返回 0。关键代码行必须能逐句翻译成“旧状态怎样变成新状态”，否则就只是模板。

\textbf{正确性。} 小于左侧最大值的位置必须参与排序，否则它前面有更大元素；大于右侧最小值的位置同理必须参与排序。

\textbf{边界实验。} 先用已经有序、逆序、重复值、局部乱序做反例检查。实验时把已经有序、逆序、重复值、局部乱序列成表，再打印关键下标和有效区间。若某个元素被写入后又被未来步骤破坏，说明区间不变量没有成立。

\textbf{复杂度变化。} 排序比较法 O(n log n)；双向扫描 O(n) 时间、O(1) 空间。

\textbf{迁移追问。} 如果题目改成“排序比较法更直观但复杂度更高”，先判断原状态是否仍然覆盖新答案集合，再决定是微调边界、改 value 语义，还是换算法。

\noindent\textbf{LeetCode 587 安装栅栏。}

\textbf{题目说明。} LeetCode 587 安装栅栏 的题面入口要先还原成具体任务：外层点构成凸包，内部点不会影响围栏边界。

\textbf{样例入口。} 拿一组点里所有边界点围成栅栏手算，内部点不会影响围栏；排序后构造下凸壳时，如果新点让边界向右转，就说明中间点在内部，要弹出。

\textbf{暴力基线。} 暴力可以检查每两个点构成的直线，若所有点都在同一侧，则这条线可能是凸包边。

\textbf{暴力过程。} 若有大量点，暴力要枚举点对并检查所有其他点。排序后从左到右构造边界，可以逐步排除内部点。

\textbf{重复工作。} 题面模型：外层点构成凸包，内部点不会影响围栏边界。暴力版的慢点要落到本题：浪费在反复判断同一批点是否在边界内。单调链维护当前凸壳，遇到右转时中间点已经不可能留在外边界。后面的优化只消掉这类重复工作，不能改变答案集合。

\textbf{优化条件。} 本题的关键观察是：排序后用单调链维护上下凸壳，遇到会让边界右转的点就弹出，题目要求保留共线边界点。这句话必须能翻译成可维护状态；如果题目条件改变后观察不再成立，就不能继续套原来的写法。

\textbf{相邻状态。} 规律是单调链用叉积维护凸壳，题目要求保留共线边界点时等号不能随便弹。把这个特例继续拆开看：先标出已经确认的前缀或后缀，再标出未知区；能被丢掉的候选，必须是以后再也不会改变已确认区域语义的候选。这样才算从例子推出数组不变量，而不是只记一次操作。

\textbf{状态复用。} 按坐标排序，分别构造下凸壳和上凸壳。用叉积判断最近两个壳点加新点是否破坏外边界；题目要求保留共线边界点时等号处理要谨慎。手算时只改被新输入影响的那一小部分，而不是回头重算完整候选。

\textbf{指针和字段。} 状态字段要能说清：hull 是当前边界点栈，cross(a, b, c) 表示转向，排序后的点序提供扫描方向。手算时按这个协议跟踪：新点让最后两条边形成右转时，倒数第二个点落到内部，弹出；若三点共线且在边界上，需要保留共线点。

\textbf{代码落点。} 关键代码是 while hull size>=2 \&\& cross(...) < 0 才弹出；若写成 <=0 会把共线边界点删掉。关键代码行必须能逐句翻译成“旧状态怎样变成新状态”，否则就只是模板。

\textbf{正确性。} 单调链每次弹出的点都位于新边形成的内部侧，不可能成为最终外边界；上下壳合并覆盖整个凸包边界。

\textbf{边界实验。} 先用所有点共线、重复点、共线边界是否保留、叉积溢出做反例检查。实验时把所有点共线、重复点、共线边界是否保留、叉积溢出列成表，再打印关键下标和有效区间。若某个元素被写入后又被未来步骤破坏，说明区间不变量没有成立。

\textbf{复杂度变化。} 暴力 O(\ensuremath{n^{3}})，单调链排序 O(n log n)、扫描 O(n)，空间 O(n)。

\textbf{迁移追问。} 如果题目改成“若只求凸包顶点不保留边上点，叉积等号处理不同”，先判断原状态是否仍然覆盖新答案集合，再决定是微调边界、改 value 语义，还是换算法。

\noindent\textbf{LeetCode 647 回文子串。}

\textbf{题目说明。} LeetCode 647 回文子串 的题面入口要先还原成具体任务：每个回文都有唯一中心或中心间隙。

\textbf{样例入口。} 拿 aaa 手算，以每个字符为中心有 3 个单字符回文，以两个字符中间为中心有 2 个 aa，以中间 a 为中心还有 aaa。

\textbf{暴力基线。} 暴力枚举所有子串，再判断每个子串是否回文，回文就计数。

\textbf{暴力过程。} 以 aaa 为例，暴力会检查 a、aa、aaa、a、aa、a。很多判断都重复比较相同中心附近的字符。

\textbf{重复工作。} 题面模型：每个回文都有唯一中心或中心间隙。暴力版的慢点要落到本题：浪费在每个子串单独判断。每个回文都有唯一中心或中心间隙，从中心扩展能一次数出围绕它的所有回文。后面的优化只消掉这类重复工作，不能改变答案集合。

\textbf{优化条件。} 本题的关键观察是：对每个中心向两边扩展，扩展成功一次就计数一次。这句话必须能翻译成可维护状态；如果题目条件改变后观察不再成立，就不能继续套原来的写法。

\textbf{相邻状态。} 规律是枚举字符中心和间隙中心，每次扩展成功就新增一个回文子串。把这个特例继续拆开看：先标出已经确认的前缀或后缀，再标出未知区；能被丢掉的候选，必须是以后再也不会改变已确认区域语义的候选。这样才算从例子推出数组不变量，而不是只记一次操作。

\textbf{状态复用。} 枚举 2n-1 个中心，每个中心向左右扩展；每扩展成功一次，说明找到一个新回文子串，count++。手算时只改被新输入影响的那一小部分，而不是回头重算完整候选。

\textbf{指针和字段。} 状态字段要能说清：left/right 是当前中心扩展边界，count 是已经找到的回文子串数。手算时按这个协议跟踪：aaa 中以中间 a 为中心可数出 a、aaa；以两个 a 的间隙为中心可数出 aa。

\textbf{代码落点。} 关键是对每个 i 调用 expand(i, i) 和 expand(i, i+1)，expand 每次字符相等就 ++count 再继续外扩。关键代码行必须能逐句翻译成“旧状态怎样变成新状态”，否则就只是模板。

\textbf{正确性。} 每个回文子串对应唯一中心；扩展时每个合法半径都会被计数一次，因此不重不漏。

\textbf{边界实验。} 先用空串、全相同字符、偶数回文、奇数回文做反例检查。实验时把空串、全相同字符、偶数回文、奇数回文列成表，再打印关键下标和有效区间。若某个元素被写入后又被未来步骤破坏，说明区间不变量没有成立。

\textbf{复杂度变化。} 暴力 O(\ensuremath{n^{3}})，中心扩展 O(\ensuremath{n^{2}}) 时间、O(1) 空间。

\textbf{迁移追问。} 如果题目改成“Manacher 算法可线性解决，但中心扩展更适合面试基础”，先判断原状态是否仍然覆盖新答案集合，再决定是微调边界、改 value 语义，还是换算法。

\noindent\textbf{LeetCode 1122 数组的相对排序。}

\textbf{题目说明。} LeetCode 1122 数组的相对排序 的题面入口要先还原成具体任务：arr2 给出部分值的自定义顺序，剩余值按自然升序。

\textbf{样例入口。} 拿 arr1=[2, 3, 1, 3, 2, 4, 6, 7, 9, 2, 19]、arr2=[2, 1, 4, 3, 9, 6] 手算，先按 arr2 顺序输出 2 的三个、1 的一个、4 的一个、3 的两个、9 和 6，剩余 7、19 升序跟在后面。

\textbf{暴力基线。} 暴力可以用自定义比较器排序 arr1：arr2 中的值按 arr2 位置排序，其他值按自然升序。

\textbf{暴力过程。} 以 arr1=[2, 3, 1, 3, 2, 4, 6, 7, 9, 2, 19]、arr2=[2, 1, 4, 3, 9, 6] 为例，先输出所有 2，再 1、4、3、9、6，剩余 7、19 升序。

\textbf{重复工作。} 题面模型：arr2 给出部分值的自定义顺序，剩余值按自然升序。暴力版的慢点要落到本题：浪费在比较排序。题目值域有限时，计数数组直接保存每个值出现次数，输出时按目标顺序消耗即可。后面的优化只消掉这类重复工作，不能改变答案集合。

\textbf{优化条件。} 本题的关键观察是：值域有限时用计数数组，先按 arr2 顺序输出计数，再按数值升序输出未出现于 arr2 的数。这句话必须能翻译成可维护状态；如果题目条件改变后观察不再成立，就不能继续套原来的写法。

\textbf{相邻状态。} 规律是值域有限时用计数表保存频次，输出顺序由 arr2 和自然升序两段共同决定。把这个特例继续拆开看：先标出已经确认的前缀或后缀，再标出未知区；能被丢掉的候选，必须是以后再也不会改变已确认区域语义的候选。这样才算从例子推出数组不变量，而不是只记一次操作。

\textbf{状态复用。} count[x] 记录 arr1 中 x 的频次。先遍历 arr2，把 count[x] 个 x 写入答案；再按数值从小到大输出剩余 count。手算时只改被新输入影响的那一小部分，而不是回头重算完整候选。

\textbf{指针和字段。} 状态字段要能说清：count 是频次表，write 是答案写入位置，arr2 决定第一段自定义顺序。手算时按这个协议跟踪：count[2]=3，所以第一段先写三个 2；arr2 消耗完后，count[7] 和 count[19] 仍未归零，按自然升序写到末尾。

\textbf{代码落点。} 关键是输出 arr2 中值后把 count[value] 清零，避免剩余阶段重复输出。关键代码行必须能逐句翻译成“旧状态怎样变成新状态”，否则就只是模板。

\textbf{正确性。} 计数表完整保存 arr1 多重集合；第一段严格按 arr2 顺序输出，第二段按值递增输出不在 arr2 中的剩余值，正好符合题意。

\textbf{边界实验。} 先用arr2 元素互不相同、arr1 中有不在 arr2 的值、重复值、值域边界、空数组做反例检查。实验时把arr2 元素互不相同、arr1 中有不在 arr2 的值、重复值、值域边界、空数组列成表，再打印关键下标和有效区间。若某个元素被写入后又被未来步骤破坏，说明区间不变量没有成立。

\textbf{复杂度变化。} 排序 O(n log n)；计数法 O(n+值域+arr2长度)，空间 O(值域)。

\textbf{迁移追问。} 如果题目改成“若值域很大，改用哈希计数加排序剩余值或自定义比较器”，先判断原状态是否仍然覆盖新答案集合，再决定是微调边界、改 value 语义，还是换算法。

\subsection{实验复盘}

追问通常会落在原地修改、稳定顺序、输出规模和整数溢出上。请准备说明为什么保存下标比保存指针更稳，为什么某个位置一旦写入就不会再被未来元素破坏。

复盘本节时先从 LeetCode 5 最长回文子串、LeetCode 26 删除有序数组重复项、LeetCode 27 移除元素 开始：每道题都先手写一个能覆盖全部候选的暴力版，再指出优化结构具体保存了哪一段历史信息，最后用相邻案例比较为什么不能互换解法。

\topic{滑动窗口与双指针}
这一节先看 LeetCode 3 无重复字符的最长子串、LeetCode 11 盛最多水的容器、LeetCode 15 三数之和 这几道 LeetCode 题，再整理同一题型的分流和推导。阅读顺序不是先背原理，而是先看案例的暴力瓶颈，再把重复出现的观察抽象成方法。

\subsection{原理和适用条件}

以 LeetCode 3 无重复字符的最长子串、LeetCode 11 盛最多水的容器、LeetCode 15 三数之和 为主线：LeetCode 3 无重复字符的最长子串：模型是“窗口保存当前无重复子串，右端每次加入一个字符。”，优化抓手是“用上次出现位置直接跳过无效左端，左边界只能向右不能回退。”；LeetCode 11 盛最多水的容器：模型是“给定每个下标的高度，选择两条竖线和横轴组成容器，返回能盛水的最大面积。”，优化抓手是“每次移动短板，因为移动长板只会缩小宽度且不会提高短板。”；LeetCode 15 三数之和：模型是“给定整数数组，返回所有不重复三元组，使三个数之和等于零。”，优化抓手是“排序后固定第一个数，剩余两数转成有序双指针；固定值和左右值都要跳过重复。”。这些案例共同推出的规律是：先枚举所有窗口，再确认右端扩展和左端收缩是否具有单调性。只有窗口状态可以增量维护，且移动左端不会漏掉未来更优答案时，滑动窗口才成立。 方法名只是复盘后的标签，真正要掌握的是从题面约束、暴力失败、状态设计到边界反例的推导链。

\begin{principlebox}{图示：从 LeetCode 案例到 滑动窗口与双指针推导链}
\begin{tabular}{@{}p{0.18\linewidth}p{0.74\linewidth}@{}}
暴力基线 & 枚举候选、完整模拟或逐个检查，先保证覆盖全部可能答案。\\
瓶颈定位 & 真正的瓶颈不是两层循环本身，而是窗口状态没有被增量维护。右端移动带来的变化只有一个新元素，左端移动移走的也只有一个旧元素，若每次重新统计，就丢掉了题目给的连续性。\\
结构选择 & 滑动窗口成立必须有单调性或固定长度边界。右端扩展让某个条件单调变好或变坏，左端收缩可以恢复合法性或缩短答案。若数组含负数、状态会来回震荡，就不能硬套窗口。\\
不变量 & 窗口不变量要写成当前窗口满足什么条件，或者当前窗口是第一个可能满足条件的最短前缀。每次移动指针之前先更新状态，移动之后状态要和区间边界一致。\\
代码落地 & C++ 中建议把窗口写成左闭右开或左闭右闭中的一种，并在全题坚持。字符计数用 \code{std::array<int, 128>} 或 26 长度数组时要说明字符集假设；泛化字符集再换哈希表。\\
\end{tabular}
\end{principlebox}

\subsection{LeetCode 案例分流}

\begin{principlebox}{从 LeetCode 案例分流：滑动窗口与双指针}
\textbf{固定长度窗口。}
代表案例：LeetCode 438 找到字符串中所有字母异位词、LeetCode 567 字符串排列。先看这些题的题面条件：窗口长度被题目直接限定或由目标串长度决定。由此推出的处理方式是：右进左出维护计数，窗口达到固定长度后检查状态。风险点：字符集、重复字符和窗口重叠是主要边界。

\textbf{可变合法窗口。}
代表案例：LeetCode 3 无重复字符最长子串、LeetCode 76 最小覆盖子串、LeetCode 209 长度最小子数组。先看这些题的题面条件：右端扩展可能破坏或满足合法性，左端可以单向收缩恢复条件。由此推出的处理方式是：维护窗口统计，满足或不满足条件时移动左端并更新答案。风险点：若数组含负数或条件不单调，普通窗口可能失效。

\textbf{有序数组双指针。}
代表案例：LeetCode 11 盛最多水容器、LeetCode 15 三数之和、LeetCode 18 四数之和、LeetCode 16 最接近三数之和。先看这些题的题面条件：数组已经有序或排序后目标关于左右指针单调。由此推出的处理方式是：根据当前和、面积或比较结果移动某一侧指针。风险点：排序会改变原下标，输出去重和 long long 溢出必须单独处理。

\textbf{窗口动态极值。}
代表案例：LeetCode 1438 绝对差不超过限制的最长连续子数组、LeetCode 239 滑动窗口最大值、LeetCode 862 和至少为 K 的最短子数组。先看这些题的题面条件：窗口合法性取决于最大值、最小值或最近 K 个状态中的最优值。由此推出的处理方式是：用单调队列或有序集合维护窗口内极值。风险点：这类题通常不再是普通计数窗口，需要证明支配关系。

\end{principlebox}

\subsection{案例矩阵}

本节案例覆盖 LeetCode 3 无重复字符的最长子串、LeetCode 11 盛最多水的容器、LeetCode 15 三数之和、LeetCode 16 最接近的三数之和、LeetCode 18 四数之和、LeetCode 76 最小覆盖子串、LeetCode 125 验证回文串、LeetCode 167 两数之和 II 输入有序数组、LeetCode 209 长度最小的子数组、LeetCode 438 找到字符串中所有字母异位词、LeetCode 567 字符串的排列、LeetCode 713 乘积小于 K 的子数组、LeetCode 862 和至少为 K 的最短子数组、LeetCode 1438 绝对差不超过限制的最长连续子数组。阅读时不要按题号背答案，而要先判断题面落在哪个分流场景：查询目标是什么，输入约束给了什么单调性或等价关系，暴力慢在哪里，优化结构保存了什么状态。

\begin{principlebox}{案例矩阵}
\textbf{LeetCode 3 无重复字符的最长子串。}
窗口保存当前无重复子串，右端每次加入一个字符。单点风险：空串、全相同字符、重复字符出现在窗口左侧之外、扩展字符集。

\textbf{LeetCode 11 盛最多水的容器。}
给定每个下标的高度，选择两条竖线和横轴组成容器，返回能盛水的最大面积。单点风险：两端高度相等、数组长度为二、面积乘法可能溢出。

\textbf{LeetCode 15 三数之和。}
给定整数数组，返回所有不重复三元组，使三个数之和等于零。单点风险：全零、重复元素很多、没有答案、结果顺序不要求但组合不能重复。

\textbf{LeetCode 16 最接近的三数之和。}
排序后固定一数，双指针寻找最接近目标的两数和。单点风险：差值相等、负数、目标很大、三数和溢出。

\textbf{LeetCode 18 四数之和。}
两层固定加一层双指针，关键是剪枝和去重。单点风险：重复元素、目标超出 int、答案很多导致输出空间主导复杂度。

\textbf{LeetCode 76 最小覆盖子串。}
给定字符串 s 和 t，在 s 中找一个最短连续子串，使它覆盖 t 的每个字符及其次数。单点风险：目标有重复字符、源串缺字符、最小窗口在开头或结尾、大小写。

\textbf{LeetCode 125 验证回文串。}
忽略非字母数字后，回文只需要两端字符向中间比较。单点风险：空串、只有标点、大小写、数字字符、Unicode 范围假设。

\textbf{LeetCode 167 两数之和 II 输入有序数组。}
有序数组中左右指针的和随指针移动单调变化。单点风险：两个数在边界、重复值、题目下标从一开始、保证有解。

\textbf{LeetCode 209 长度最小的子数组。}
给定正整数数组和目标值，返回和至少为目标值的最短连续子数组长度。单点风险：不存在答案、单元素达到目标、全正是必要条件。

\textbf{LeetCode 438 找到字符串中所有字母异位词。}
固定长度窗口内字符频次等于目标频次即为答案。单点风险：目标比源串长、重复字符、字符集固定、答案重叠。

\textbf{LeetCode 567 字符串的排列。}
目标串任意排列出现在源串中，等价于固定长度窗口频次相同。单点风险：目标比源串长、重复字符、字符集固定、窗口重叠。

\textbf{LeetCode 713 乘积小于 K 的子数组。}
正整数数组的窗口乘积随右端扩展变大、随左端收缩变小。单点风险：k 小于等于一、元素都为一、乘积溢出、窗口缩到空、严格小于而不是小于等于。

\textbf{LeetCode 862 和至少为 K 的最短子数组。}
数组允许负数，普通滑动窗口的和不再随左端移动单调变化。单点风险：负数、答案不存在、K 为负或零、队尾支配关系。

\textbf{LeetCode 1438 绝对差不超过限制的最长连续子数组。}
窗口合法性取决于窗口内最大值和最小值之差。单点风险：limit 为零、重复元素、最大最小同时过期、窗口收缩多次。

\end{principlebox}

\subsection{共性落点}

\begin{principlebox}{本组共性落点}
\textbf{慢版本。} 暴力做法枚举所有左端和右端，再检查窗口是否合法。对于字符串和正数数组，这种写法会把相邻窗口中大量相同内容反复计算。
\textbf{瓶颈。} 真正的瓶颈不是两层循环本身，而是窗口状态没有被增量维护。右端移动带来的变化只有一个新元素，左端移动移走的也只有一个旧元素，若每次重新统计，就丢掉了题目给的连续性。
\textbf{状态字段。} 左端、右端、窗口计数或当前双指针和、合法性指标、当前答案；每次移动边界后，状态必须和候选区间或窗口内容一致。
\textbf{不变量。} 窗口不变量要写成当前窗口满足什么条件，或者当前窗口是第一个可能满足条件的最短前缀。每次移动指针之前先更新状态，移动之后状态要和区间边界一致。
\textbf{实现检查。} 先更新右端进入，再判断合法性，收缩时先记录答案还是先移走左端要固定；不要让左端回退。
\textbf{C++ 落地。} C++ 中建议把窗口写成左闭右开或左闭右闭中的一种，并在全题坚持。字符计数用 \code{std::array<int, 128>} 或 26 长度数组时要说明字符集假设；泛化字符集再换哈希表。
\textbf{证明抓手。} 证明从左端淘汰开始：被移走的左端对应窗口已经检查完，或继续保留只会让状态更差。
\end{principlebox}

\subsection{案例推导}

\noindent\textbf{LeetCode 3 无重复字符的最长子串。}

\textbf{题目说明。} LeetCode 3 无重复字符的最长子串 的题面入口要先还原成具体任务：窗口保存当前无重复子串，右端每次加入一个字符。

\textbf{样例入口。} 拿字符串 abca 手算。窗口 abc 合法，右端加入第二个 a 后，真正冲突的是上一次 a 的位置；左端从 0 跳到 1 后，bca 重新合法。再试 abba：处理第二个 b 时左端跳到 2，后面遇到 a 的上次位置在窗口左侧之外，不能让 left 回退。

\textbf{暴力基线。} 暴力固定每个起点 left，向右逐个加入字符，用集合检查是否重复；一旦重复就停止当前起点，换下一个起点重新建集合。

\textbf{暴力过程。} 以 s=abcabcbb 为例，暴力会检查 a、ab、abc，到 abca 时发现 a 重复；再从 b 重新检查 b、bc、bca。相邻起点的窗口 b、c 其实已经在上一轮出现过。

\textbf{重复工作。} 题面模型：窗口保存当前无重复子串，右端每次加入一个字符。暴力版的慢点要落到本题：浪费在每换一个 left 就重新统计一段已经看过的字符。真正变化只是左端离开若干旧字符，右端继续加入一个新字符。后面的优化只消掉这类重复工作，不能改变答案集合。

\textbf{优化条件。} 本题的关键观察是：用上次出现位置直接跳过无效左端，左边界只能向右不能回退。这句话必须能翻译成可维护状态；如果题目条件改变后观察不再成立，就不能继续套原来的写法。

\textbf{相邻状态。} 规律是左边界只向右，状态保存每个字符最近出现位置。把这个特例继续拆开看：右端进入只带来一个新元素，左端离开只删掉一个旧元素；只有当旧左端被移走后不可能再产生更优答案时，窗口才可以单向推进。若这个单调淘汰条件不存在，就要换成前缀和、队列或其他状态。

\textbf{状态复用。} 保存窗口左端 left、右端 right、每个字符最近出现位置 last、当前最长答案 best。遇到重复字符时，left 直接跳到 last[ch]+1，但只能向右跳。手算时只改被新输入影响的那一小部分，而不是回头重算完整候选。

\textbf{指针和字段。} 状态字段要能说清：right 是当前加入字符位置，left 是当前无重复窗口起点，last[ch] 是字符 ch 上一次出现下标，best 是已经见过的最大合法窗口长度。手算时按这个协议跟踪：手算 abba：right=0 加 a，left=0；right=1 加 b；right=2 再遇 b，left 跳到 2；right=3 遇 a，但上次 a 在 0，已经在窗口外，left 不能回退。

\textbf{代码落点。} 关键代码是 left=max(left, last[ch]+1)，然后 last[ch]=right，最后用 right-left+1 更新答案。顺序不能反，否则会把当前字符当成自己的历史位置。关键代码行必须能逐句翻译成“旧状态怎样变成新状态”，否则就只是模板。

\textbf{正确性。} left 跳过的是所有仍包含重复 ch 的起点，这些起点不可能形成以 right 结尾的合法窗口；left 不回退保证每个字符至多进出窗口一次。

\textbf{边界实验。} 先用空串、全相同字符、重复字符出现在窗口左侧之外、扩展字符集做反例检查。实验时重点跑空串、全相同字符、重复字符出现在窗口左侧之外、扩展字符集，并打印左右端、窗口计数和答案。若左端发生回退，或者窗口失去合法性后仍更新答案，就能很快定位错误。

\textbf{复杂度变化。} 暴力最坏 O(\ensuremath{n^{2}})，优化后每个 right 扫一次，哈希或数组访问均摊 O(1)，总时间 O(n)，空间取决于字符集。

\textbf{迁移追问。} 如果题目改成“若要求恰好包含 k 种字符，窗口条件从无重复变成种类计数”，先判断原状态是否仍然覆盖新答案集合，再决定是微调边界、改 value 语义，还是换算法。

\noindent\textbf{LeetCode 11 盛最多水的容器。}

\textbf{题目说明。} LeetCode 11 盛最多水的容器给定数组 height，height[i] 表示第 i 条竖线高度。选择两条线与 x 轴形成容器，返回能容纳的最大水量。

\textbf{样例入口。} 样例 height=[1, 8, 6, 2, 5, 4, 8, 3, 7]，输出 49；选择下标 1 和 8，宽度 7，短板高度 7，面积 49。

\textbf{暴力基线。} 暴力枚举所有 left<right 的下标对，面积等于 (right-left)*min(height[left], height[right])，取最大值。

\textbf{暴力过程。} 以 [1, 8, 6, 2, 5, 4, 8, 3, 7] 为例，暴力会试 (0, 1)、(0, 2) 一直到 (0, 8)，再试 (1, 2)、(1, 3)。大量候选只是宽度和短板变化。

\textbf{重复工作。} 题面模型：给定每个下标的高度，选择两条竖线和横轴组成容器，返回能盛水的最大面积。暴力版的慢点要落到本题：浪费在逐个试所有右端。若当前短板是 left，保留 left 再把 right 左移只会让宽度变小，短板高度不可能超过 height[left]，这些候选都被当前状态支配。后面的优化只消掉这类重复工作，不能改变答案集合。

\textbf{优化条件。} 本题的关键观察是：每次移动短板，因为移动长板只会缩小宽度且不会提高短板。这句话必须能翻译成可维护状态；如果题目条件改变后观察不再成立，就不能继续套原来的写法。

\textbf{相邻状态。} 规律是每轮淘汰短板那一侧，因为所有保留它且宽度更小的候选都被当前面积上界支配。把这个特例继续拆开看：右端进入只带来一个新元素，左端离开只删掉一个旧元素；只有当旧左端被移走后不可能再产生更优答案时，窗口才可以单向推进。若这个单调淘汰条件不存在，就要换成前缀和、队列或其他状态。

\textbf{状态复用。} 双指针从两端开始，维护 left、right 和 best。每轮计算当前面积，然后移动高度较小的一侧。手算时只改被新输入影响的那一小部分，而不是回头重算完整候选。

\textbf{指针和字段。} 状态字段要能说清：left/right 是当前最宽候选的两端，shorter 是限制高度的一侧，best 是已见过最大面积。手算时按这个协议跟踪：样例中 left=0 高 1、right=8 高 7，面积 8。短板是 left，移动 left 到 1；此时高 8 和 7，面积 49，更新答案。

\textbf{代码落点。} 关键代码是 if (height[left] < height[right]) ++left; else --right; 。移动长板无法提高短板上界，所以不移动长板。关键代码行必须能逐句翻译成“旧状态怎样变成新状态”，否则就只是模板。

\textbf{正确性。} 每次移动短板时，所有包含这个短板且宽度更小的候选面积都不可能超过当前面积上界，因此可以一次排除一批候选。

\textbf{边界实验。} 先用两端高度相等、数组长度为二、面积乘法可能溢出做反例检查。实验时重点跑两端高度相等、数组长度为二、面积乘法可能溢出，并打印左右端、窗口计数和答案。若左端发生回退，或者窗口失去合法性后仍更新答案，就能很快定位错误。

\textbf{复杂度变化。} 暴力 O(\ensuremath{n^{2}})，双指针每轮移动一端，总共 O(n) 轮，时间 O(n)，空间 O(1)。

\textbf{迁移追问。} 如果题目改成“接雨水计算每个位置贡献，不能把本题双指针证明直接搬过去”，先判断原状态是否仍然覆盖新答案集合，再决定是微调边界、改 value 语义，还是换算法。

\noindent\textbf{LeetCode 15 三数之和。}

\textbf{题目说明。} LeetCode 15 三数之和给定整数数组 nums，要求找出所有不重复三元组 [a, b, c]，使 a+b+c=0。

\textbf{样例入口。} 样例 nums=[-1, 0, 1, 2, -1, -4]，输出 [[-1, -1, 2], [-1, 0, 1]]；结果里不能出现重复三元组。

\textbf{暴力基线。} 暴力三层循环枚举 i、j、k，检查 nums[i]+nums[j]+nums[k] 是否为 0，再用集合去掉重复三元组。

\textbf{暴力过程。} 以 [-1, 0, 1, 2, -1, -4] 为例，暴力会枚举 (-1, 0, 1)、(-1, 0, 2)、(-1, 0, -1) 等所有三元组；两个 -1 会制造重复答案。

\textbf{重复工作。} 题面模型：给定整数数组，返回所有不重复三元组，使三个数之和等于零。暴力版的慢点要落到本题：浪费在第三层枚举和事后去重。排序后固定一个数，剩余两个数变成有序两数之和，可以根据当前和一次排除一侧。后面的优化只消掉这类重复工作，不能改变答案集合。

\textbf{优化条件。} 本题的关键观察是：排序后固定第一个数，剩余两数转成有序双指针；固定值和左右值都要跳过重复。这句话必须能翻译成可维护状态；如果题目条件改变后观察不再成立，就不能继续套原来的写法。

\textbf{相邻状态。} 规律是排序把三数问题拆成枚举固定点加双指针，并在每层处理去重。把这个特例继续拆开看：右端进入只带来一个新元素，左端离开只删掉一个旧元素；只有当旧左端被移走后不可能再产生更优答案时，窗口才可以单向推进。若这个单调淘汰条件不存在，就要换成前缀和、队列或其他状态。

\textbf{状态复用。} 先排序。外层 i 固定第一个数，内层 left=i+1、right=n-1 搜索 target=-nums[i]；找到答案后跳过相同 left/right 值。手算时只改被新输入影响的那一小部分，而不是回头重算完整候选。

\textbf{指针和字段。} 状态字段要能说清：i 是本层固定点，left/right 是右侧有序区间的候选两端，sum 是三数和，res 保存已去重结果。手算时按这个协议跟踪：排序后 [-4, -1, -1, 0, 1, 2]。i=1 固定 -1，left 指向 -1，right 指向 2，和为 0，记录 [-1, -1, 2]，然后跳过重复。

\textbf{代码落点。} 关键代码是 sum<0 时 ++left，sum>0 时 --right；i>0 且 nums[i]==nums[i-1] 要 continue，命中后 left/right 也要跳过重复值。关键代码行必须能逐句翻译成“旧状态怎样变成新状态”，否则就只是模板。

\textbf{正确性。} 排序后，当前和太小，固定 i 和 right 时更小的 left 都不可能达标；当前和太大，固定 i 和 left 时更大的 right 都不可能达标。去重发生在同一层，避免漏掉不同层组合。

\textbf{边界实验。} 先用全零、重复元素很多、没有答案、结果顺序不要求但组合不能重复做反例检查。实验时重点跑全零、重复元素很多、没有答案、结果顺序不要求但组合不能重复，并打印左右端、窗口计数和答案。若左端发生回退，或者窗口失去合法性后仍更新答案，就能很快定位错误。

\textbf{复杂度变化。} 暴力 O(\ensuremath{n^{3}}) 加去重开销；排序 O(n log n)，外层 n 次、内层双指针线性，总时间 O(\ensuremath{n^{2}})，输出空间另计。

\textbf{迁移追问。} 如果题目改成“四数之和再外层固定一个数，并用 long long 防止目标和溢出”，先判断原状态是否仍然覆盖新答案集合，再决定是微调边界、改 value 语义，还是换算法。

\noindent\textbf{LeetCode 16 最接近的三数之和。}

\textbf{题目说明。} LeetCode 16 最接近的三数之和 的题面入口要先还原成具体任务：排序后固定一数，双指针寻找最接近目标的两数和。

\textbf{样例入口。} 拿 [-1, 2, 1, -4]、target=1 手算，排序后固定 -1，左右指针在 1 和 2 上得到和 2，距离目标只差 1；若和偏大，只能让右指针左移，因为右侧数再大只会更远。

\textbf{暴力基线。} 暴力枚举所有三元组，计算三数和与 target 的绝对差，保留差值最小的和。

\textbf{暴力过程。} 以 nums=[-1, 2, 1, -4]、target=1 为例，暴力会试 (-1, 2, 1)、(-1, 2, -4)、(-1, 1, -4) 等全部组合。

\textbf{重复工作。} 题面模型：排序后固定一数，双指针寻找最接近目标的两数和。暴力版的慢点要落到本题：浪费在固定一个数后仍然枚举剩余两数。排序后当前和偏小或偏大能决定移动哪一端，不需要再试被单调性排除的一侧。后面的优化只消掉这类重复工作，不能改变答案集合。

\textbf{优化条件。} 本题的关键观察是：每次根据当前和与目标的大小移动指针，同时更新最小绝对差。这句话必须能翻译成可维护状态；如果题目条件改变后观察不再成立，就不能继续套原来的写法。

\textbf{相邻状态。} 规律是每次先更新最小差，再按当前和与目标的大小排除一侧候选；差值相等和三数和溢出要单独检查。把这个特例继续拆开看：右端进入只带来一个新元素，左端离开只删掉一个旧元素；只有当旧左端被移走后不可能再产生更优答案时，窗口才可以单向推进。若这个单调淘汰条件不存在，就要换成前缀和、队列或其他状态。

\textbf{状态复用。} 排序后固定 i，用 left/right 夹逼。每轮先用当前 sum 更新 best\_sum，再根据 sum 与 target 的大小移动 left 或 right。手算时只改被新输入影响的那一小部分，而不是回头重算完整候选。

\textbf{指针和字段。} 状态字段要能说清：best\_sum 是当前最接近目标的三数和，best\_diff 是差值；left/right 仍表示固定 i 后的两数候选。手算时按这个协议跟踪：排序为 [-4, -1, 1, 2]。固定 -4 时和偏小，left 右移；固定 -1 时  -1+1+2=2，距离 target 为 1，更新 best。

\textbf{代码落点。} 关键代码是 if (abs(sum-target)<abs(best-target)) best=sum; 然后 sum<target 时 ++left，否则 --right；sum==target 可直接返回。关键代码行必须能逐句翻译成“旧状态怎样变成新状态”，否则就只是模板。

\textbf{正确性。} 在固定 i 下，sum 小于目标时右端再左移只会更小，所以只能左端右移；sum 大于目标时同理只能右端左移。

\textbf{边界实验。} 先用差值相等、负数、目标很大、三数和溢出做反例检查。实验时重点跑差值相等、负数、目标很大、三数和溢出，并打印左右端、窗口计数和答案。若左端发生回退，或者窗口失去合法性后仍更新答案，就能很快定位错误。

\textbf{复杂度变化。} 暴力 O(\ensuremath{n^{3}})，排序加双指针 O(\ensuremath{n^{2}})，额外空间 O(1) 或取决于排序实现。

\textbf{迁移追问。} 如果题目改成“若要求返回所有最接近组合，需要记录同差值并去重”，先判断原状态是否仍然覆盖新答案集合，再决定是微调边界、改 value 语义，还是换算法。

\noindent\textbf{LeetCode 18 四数之和。}

\textbf{题目说明。} LeetCode 18 四数之和 的题面入口要先还原成具体任务：两层固定加一层双指针，关键是剪枝和去重。

\textbf{样例入口。} 拿 [1, 0, -1, 0, -2, 2]、target=0 手算，排序后先固定 -2，再固定 -1，剩余两数用双指针找 3。找到 [-2, -1, 1, 2] 后，左右指针相同值必须跳过，否则会重复输出。

\textbf{暴力基线。} 暴力四层循环枚举四个下标，检查和是否等于 target，再用集合去重。

\textbf{暴力过程。} 以 [1, 0, -1, 0, -2, 2]、target=0 为例，暴力会试所有四元组，[-2, -1, 1, 2] 和含重复 0 的组合都要事后去重。

\textbf{重复工作。} 题面模型：两层固定加一层双指针，关键是剪枝和去重。暴力版的慢点要落到本题：浪费在后两层仍然盲枚举。排序后两层固定，剩余两数是有序配对问题；同层重复值可以直接跳过。后面的优化只消掉这类重复工作，不能改变答案集合。

\textbf{优化条件。} 本题的关键观察是：排序提供单调性，外层可用最小可能和和最大可能和提前跳过。这句话必须能翻译成可维护状态；如果题目条件改变后观察不再成立，就不能继续套原来的写法。

\textbf{相邻状态。} 规律是两层固定把四数降成有序两数，外层用最小可能和与最大可能和剪枝，内外层都要去重。把这个特例继续拆开看：右端进入只带来一个新元素，左端离开只删掉一个旧元素；只有当旧左端被移走后不可能再产生更优答案时，窗口才可以单向推进。若这个单调淘汰条件不存在，就要换成前缀和、队列或其他状态。

\textbf{状态复用。} 排序后固定 i 和 j，left=j+1、right=n-1 搜索 target-nums[i]-nums[j]。外层可用最小可能和、最大可能和剪枝。手算时只改被新输入影响的那一小部分，而不是回头重算完整候选。

\textbf{指针和字段。} 状态字段要能说清：i/j 是前两个固定值，left/right 是剩余区间候选，sum 用 long long 保存，res 保存去重四元组。手算时按这个协议跟踪：排序后 [-2, -1, 0, 0, 1, 2]。i=-2、j=-1 时，left=0、right=2，四数和为 -1，left 右移；随后找到 [-2, -1, 1, 2]。

\textbf{代码落点。} 关键代码和三数之和一样，但外层有两层去重。sum 必须用 long long，避免 target 和 nums 值接近 int 边界时溢出。关键代码行必须能逐句翻译成“旧状态怎样变成新状态”，否则就只是模板。

\textbf{正确性。} 两层固定后，有序双指针的排除理由和两数之和相同；同层去重只跳过值相同且角色相同的候选，不会跳过需要使用的重复元素。

\textbf{边界实验。} 先用重复元素、目标超出 int、答案很多导致输出空间主导复杂度做反例检查。实验时重点跑重复元素、目标超出 int、答案很多导致输出空间主导复杂度，并打印左右端、窗口计数和答案。若左端发生回退，或者窗口失去合法性后仍更新答案，就能很快定位错误。

\textbf{复杂度变化。} 暴力 O(\ensuremath{n^{4}})，排序后 O(\ensuremath{n^{3}})，输出空间另计。

\textbf{迁移追问。} 如果题目改成“k 数之和可以递归降维，但工程实现要控制重复和边界”，先判断原状态是否仍然覆盖新答案集合，再决定是微调边界、改 value 语义，还是换算法。

\noindent\textbf{LeetCode 76 最小覆盖子串。}

\textbf{题目说明。} LeetCode 76 最小覆盖子串给定字符串 s 和 t，要求在 s 中找最短连续子串，使它包含 t 中每个字符及其次数；不存在则返回空串。

\textbf{样例入口。} 样例 s=ADOBECODEBANC，t=ABC，输出 BANC；它是覆盖 A、B、C 的最短连续子串。

\textbf{暴力基线。} 暴力枚举 s 的每个子串，重新统计字符频次，检查是否覆盖 t 的每个字符及其次数。

\textbf{暴力过程。} 以 s=ADOBECODEBANC、t=ABC 为例，暴力会检查 A、AD、ADO...直到 ADOBEC 才第一次覆盖；下一轮从 D 开始又重新统计 DOBEC。

\textbf{重复工作。} 题面模型：给定字符串 s 和 t，在 s 中找一个最短连续子串，使它覆盖 t 的每个字符及其次数。暴力版的慢点要落到本题：浪费在相邻子串重复统计字符。右端扩展只新增一个字符，左端收缩只移走一个字符。后面的优化只消掉这类重复工作，不能改变答案集合。

\textbf{优化条件。} 本题的关键观察是：窗口内字符计数必须覆盖目标计数，满足后尽量收缩左端；用 formed 记录满足需求的字符种类，避免每次全量比较。这句话必须能翻译成可维护状态；如果题目条件改变后观察不再成立，就不能继续套原来的写法。

\textbf{相邻状态。} 规律是右端负责获得可行，左端负责在可行时尽量缩短；计数表要区分已有字符和需求字符。把这个特例继续拆开看：右端进入只带来一个新元素，左端离开只删掉一个旧元素；只有当旧左端被移走后不可能再产生更优答案时，窗口才可以单向推进。若这个单调淘汰条件不存在，就要换成前缀和、队列或其他状态。

\textbf{状态复用。} need 保存目标频次，window 保存当前窗口频次，formed 或 required\_count 记录当前满足需求的字符种类；满足后尽量收缩 left。手算时只改被新输入影响的那一小部分，而不是回头重算完整候选。

\textbf{指针和字段。} 状态字段要能说清：right 负责扩展获得覆盖，left 负责在覆盖时缩短窗口，best\_left/best\_len 记录最短答案。手算时按这个协议跟踪：ADOBEC 第一次覆盖 ABC，记录长度 6；左端移走 A 后不覆盖，停止收缩。继续扩展到 BANC 时再次覆盖，收缩得到长度 4。

\textbf{代码落点。} 关键顺序是加入 s[right] 后更新 formed；while formed==required 时先更新答案，再移走 s[left] 并可能降低 formed。关键代码行必须能逐句翻译成“旧状态怎样变成新状态”，否则就只是模板。

\textbf{正确性。} 每个 right 下，while 循环会把 left 推到刚好不能再缩的位置，因此不会错过以该 right 结尾的最短覆盖窗口。

\textbf{边界实验。} 先用目标有重复字符、源串缺字符、最小窗口在开头或结尾、大小写做反例检查。实验时重点跑目标有重复字符、源串缺字符、最小窗口在开头或结尾、大小写，并打印左右端、窗口计数和答案。若左端发生回退，或者窗口失去合法性后仍更新答案，就能很快定位错误。

\textbf{复杂度变化。} 暴力至少 O(\ensuremath{n^{2}}*字符集)，滑动窗口每个字符进出一次，时间 O(n+字符集)，空间 O(字符集)。

\textbf{迁移追问。} 如果题目改成“若字符集固定，数组计数更快；若是 Unicode，哈希表更通用”，先判断原状态是否仍然覆盖新答案集合，再决定是微调边界、改 value 语义，还是换算法。

\noindent\textbf{LeetCode 125 验证回文串。}

\textbf{题目说明。} LeetCode 125 验证回文串 的题面入口要先还原成具体任务：忽略非字母数字后，回文只需要两端字符向中间比较。

\textbf{样例入口。} 拿 A man, a plan, a canal: Panama 手算，忽略空格和标点后得到 amanaplanacanalpanama，左右两端逐个比较都相同。

\textbf{暴力基线。} 暴力可以先构造一个只含字母数字且统一小写的新字符串，再反转或双指针比较。

\textbf{暴力过程。} 以 A man, a plan, a canal: Panama 为例，过滤后是 amanaplanacanalpanama，再判断它是否回文。

\textbf{重复工作。} 题面模型：忽略非字母数字后，回文只需要两端字符向中间比较。暴力版的慢点要落到本题：浪费在新建过滤字符串。原字符串两端向中间走时，可以边跳过无关字符边比较。后面的优化只消掉这类重复工作，不能改变答案集合。

\textbf{优化条件。} 本题的关键观察是：左右指针分别跳过非法字符，再比较小写化后的字符，相等则同时收缩。这句话必须能翻译成可维护状态；如果题目条件改变后观察不再成立，就不能继续套原来的写法。

\textbf{相邻状态。} 规律是左右指针跳过非字母数字字符，比较时统一大小写，指针只向中间移动。把这个特例继续拆开看：右端进入只带来一个新元素，左端离开只删掉一个旧元素；只有当旧左端被移走后不可能再产生更优答案时，窗口才可以单向推进。若这个单调淘汰条件不存在，就要换成前缀和、队列或其他状态。

\textbf{状态复用。} left/right 从两端开始，跳过非字母数字字符，比较 tolower 后的字符；相等则同时收缩。手算时只改被新输入影响的那一小部分，而不是回头重算完整候选。

\textbf{指针和字段。} 状态字段要能说清：left/right 是尚未确认的两端字符位置，合法字符判断决定是否参与比较。手算时按这个协议跟踪：左端 A 和右端 a 统一小写后相等；中间遇到空格、逗号、冒号都跳过。

\textbf{代码落点。} 关键是调用 isalnum 和 tolower 时用 unsigned char 转换，避免 char 为负时行为不稳；比较失败立即返回 false。关键代码行必须能逐句翻译成“旧状态怎样变成新状态”，否则就只是模板。

\textbf{正确性。} 回文只要求过滤后的第 k 个字符等于倒数第 k 个字符；双指针跳过无关字符后比较的正是这两者。

\textbf{边界实验。} 先用空串、只有标点、大小写、数字字符、Unicode 范围假设做反例检查。实验时重点跑空串、只有标点、大小写、数字字符、Unicode 范围假设，并打印左右端、窗口计数和答案。若左端发生回退，或者窗口失去合法性后仍更新答案，就能很快定位错误。

\textbf{复杂度变化。} 构造法 O(n) 时间和 O(n) 空间；原地双指针 O(n) 时间、O(1) 空间。

\textbf{迁移追问。} 如果题目改成“若要求最多删除一个字符，可以在第一次失配时分支检查两个子区间”，先判断原状态是否仍然覆盖新答案集合，再决定是微调边界、改 value 语义，还是换算法。

\noindent\textbf{LeetCode 167 两数之和 II 输入有序数组。}

\textbf{题目说明。} LeetCode 167 两数之和 II 输入有序数组 的题面入口要先还原成具体任务：有序数组中左右指针的和随指针移动单调变化。

\textbf{样例入口。} 拿 [2, 7, 11, 15]、target=9 手算，左右和为 17 太大，右指针左移；2+7 正好等于 9。

\textbf{暴力基线。} 暴力枚举所有 i<j 的下标对，检查 numbers[i]+numbers[j] 是否等于 target。

\textbf{暴力过程。} 以 [2, 7, 11, 15]、target=9 为例，暴力会试 2+7、2+11、2+15...；但数组有序，2+15 太大时，右端继续更大没有意义。

\textbf{重复工作。} 题面模型：有序数组中左右指针的和随指针移动单调变化。暴力版的慢点要落到本题：浪费在没有利用有序性。当前和太小，固定右端时更小的左端都不可能达标；当前和太大，固定左端时更大的右端都不可能达标。后面的优化只消掉这类重复工作，不能改变答案集合。

\textbf{优化条件。} 本题的关键观察是：当前和太小左指针右移，太大右指针左移，等于目标返回。这句话必须能翻译成可维护状态；如果题目条件改变后观察不再成立，就不能继续套原来的写法。

\textbf{相邻状态。} 规律是有序数组让和随左指针右移变大、随右指针左移变小，因此每次能排除一侧候选。把这个特例继续拆开看：右端进入只带来一个新元素，左端离开只删掉一个旧元素；只有当旧左端被移走后不可能再产生更优答案时，窗口才可以单向推进。若这个单调淘汰条件不存在，就要换成前缀和、队列或其他状态。

\textbf{状态复用。} left 从 0 开始，right 从 n-1 开始。每轮计算 sum，sum 小则 left++，sum 大则 right--，相等返回 1 基下标。手算时只改被新输入影响的那一小部分，而不是回头重算完整候选。

\textbf{指针和字段。} 状态字段要能说清：left/right 是当前候选对，sum 是当前两数和，target 决定排除哪一侧。手算时按这个协议跟踪：[2, 7, 11, 15] 中 2+15=17 太大，right 左移；2+11=13 仍大，再左移；2+7=9 返回 [1, 2]。

\textbf{代码落点。} 关键代码是 sum<target 时 ++left，sum>target 时 --right。返回时要加 1，因为题目要求 1 基下标。关键代码行必须能逐句翻译成“旧状态怎样变成新状态”，否则就只是模板。

\textbf{正确性。} 有序性保证移动一侧能排除一整批不可能候选，且剩余区间仍包含可能答案。

\textbf{边界实验。} 先用两个数在边界、重复值、题目下标从一开始、保证有解做反例检查。实验时重点跑两个数在边界、重复值、题目下标从一开始、保证有解，并打印左右端、窗口计数和答案。若左端发生回退，或者窗口失去合法性后仍更新答案，就能很快定位错误。

\textbf{复杂度变化。} 暴力 O(\ensuremath{n^{2}})，双指针 O(n) 时间、O(1) 空间。

\textbf{迁移追问。} 如果题目改成“无序版本通常用哈希表，不能直接套双指针”，先判断原状态是否仍然覆盖新答案集合，再决定是微调边界、改 value 语义，还是换算法。

\noindent\textbf{LeetCode 209 长度最小的子数组。}

\textbf{题目说明。} LeetCode 209 长度最小的子数组给定正整数数组 nums 和目标 target，返回和至少为 target 的最短连续子数组长度；不存在则返回 0。

\textbf{样例入口。} 样例 target=7，nums=[2, 3, 1, 2, 4, 3]，输出 2；最短合法连续段是 [4, 3]。

\textbf{暴力基线。} 暴力枚举每个起点 left，向右累计 sum，第一次 sum>=target 时更新长度，再换下一个起点。

\textbf{暴力过程。} 以 target=7、nums=[2, 3, 1, 2, 4, 3] 为例，从 0 开始要扩到 [2, 3, 1, 2]；从 1 开始又重新累加 [3, 1, 2, 4]。

\textbf{重复工作。} 题面模型：给定正整数数组和目标值，返回和至少为目标值的最短连续子数组长度。暴力版的慢点要落到本题：浪费在相邻起点重复累加。因为 nums 全为正数，右端增加只会让 sum 变大，左端移走只会让 sum 变小。后面的优化只消掉这类重复工作，不能改变答案集合。

\textbf{优化条件。} 本题的关键观察是：正数数组让窗口和随右端增加、随左端减少；右端扩展到满足目标后，不断左移缩短并更新答案。这句话必须能翻译成可维护状态；如果题目条件改变后观察不再成立，就不能继续套原来的写法。

\textbf{相邻状态。} 规律是全正数保证右扩和变大、左缩和变小，所以每个左端只需被移动一次。把这个特例继续拆开看：右端进入只带来一个新元素，左端离开只删掉一个旧元素；只有当旧左端被移走后不可能再产生更优答案时，窗口才可以单向推进。若这个单调淘汰条件不存在，就要换成前缀和、队列或其他状态。

\textbf{状态复用。} left、right 维护当前窗口和 sum。right 每次加入一个数；当 sum>=target 时，循环收缩 left 并更新最短长度。手算时只改被新输入影响的那一小部分，而不是回头重算完整候选。

\textbf{指针和字段。} 状态字段要能说清：sum 是 nums[left..right] 的和，answer 是已见过最短合法长度；left 只在窗口合法时右移。手算时按这个协议跟踪：样例中 right 到下标 3 时 sum=8，记录长度 4，移走 2 后 sum=6 停止；后来窗口 [4, 3] sum=7，答案更新为 2。

\textbf{代码落点。} 关键代码是 sum += nums[right]; while (sum >= target) \{ ans=min(ans, right-left+1); sum-=nums[left++]; \}。关键代码行必须能逐句翻译成“旧状态怎样变成新状态”，否则就只是模板。

\textbf{正确性。} 正数保证收缩 left 后 sum 单调下降；对每个 right，while 会找到以它结尾的所有可缩短合法窗口，不会漏最短值。

\textbf{边界实验。} 先用不存在答案、单元素达到目标、全正是必要条件做反例检查。实验时重点跑不存在答案、单元素达到目标、全正是必要条件，并打印左右端、窗口计数和答案。若左端发生回退，或者窗口失去合法性后仍更新答案，就能很快定位错误。

\textbf{复杂度变化。} 暴力 O(\ensuremath{n^{2}})，窗口中每个元素至多进入一次、离开一次，时间 O(n)，空间 O(1)。

\textbf{迁移追问。} 如果题目改成“若含负数，需要前缀和加单调队列”，先判断原状态是否仍然覆盖新答案集合，再决定是微调边界、改 value 语义，还是换算法。

\noindent\textbf{LeetCode 438 找到字符串中所有字母异位词。}

\textbf{题目说明。} LeetCode 438 找到字符串中所有字母异位词 的题面入口要先还原成具体任务：固定长度窗口内字符频次等于目标频次即为答案。

\textbf{样例入口。} 拿 s=cbaebabacd、p=abc 手算，长度为 3 的窗口 cba 频次和 abc 相同，起点 0 是答案；窗口右移时只删除 c、加入 e。

\textbf{暴力基线。} 先统计 p 的频次 target，再枚举 s 中每个长度 k=p.size() 的窗口，为每个窗口重新统计 window 后比较。

\textbf{暴力过程。} 以 s=cbaebabacd、p=abc 为例，暴力窗口依次是 cba、bae、aeb、eba、bab、aba、bac、acd。每个窗口都重新数 k 个字符。

\textbf{重复工作。} 题面模型：固定长度窗口内字符频次等于目标频次即为答案。暴力版的慢点要落到本题：相邻窗口只差一个出去字符和一个进来字符。从 cba 到 bae，只有 c 出去、e 进来，b 和 a 不该重数。后面的优化只消掉这类重复工作，不能改变答案集合。

\textbf{优化条件。} 本题的关键观察是：右进左出维护计数，窗口长度达到目标长度后检查差异。这句话必须能翻译成可维护状态；如果题目条件改变后观察不再成立，就不能继续套原来的写法。

\textbf{相邻状态。} 规律是固定长度窗口维护字符频次差，重叠答案不能漏。把这个特例继续拆开看：右端进入只带来一个新元素，左端离开只删掉一个旧元素；只有当旧左端被移走后不可能再产生更优答案时，窗口才可以单向推进。若这个单调淘汰条件不存在，就要换成前缀和、队列或其他状态。

\textbf{状态复用。} k 来自 p.length()。先统计第一个窗口，之后每次窗口右移一格：window[s[i-1]]--，window[s[i+k-1]]++。手算时只改被新输入影响的那一小部分，而不是回头重算完整候选。

\textbf{指针和字段。} 状态字段要能说清：i 是当前窗口起点，out=s[i-1] 是离开字符，in=s[i+k-1] 是进入字符，target/window 是 26 长度频次数组。手算时按这个协议跟踪：第一个窗口 cba 与 abc 频次相同，记录 0；滑到 bae 后 c 计数减一、e 加一；滑到 bac 时频次再次相同，记录 6。

\textbf{代码落点。} 关键代码是固定窗口右进左出。若用 unordered\_map，计数减到 0 的 key 要 erase；若题目限定小写字母，用 array/vector 长度 26 更简单。关键代码行必须能逐句翻译成“旧状态怎样变成新状态”，否则就只是模板。

\textbf{正确性。} 每个被检查窗口长度都等于 k，window 频次始终精确表示当前窗口；频次等于 target 当且仅当该窗口是 p 的异位词。

\textbf{边界实验。} 先用目标比源串长、重复字符、字符集固定、答案重叠做反例检查。实验时重点跑目标比源串长、重复字符、字符集固定、答案重叠，并打印左右端、窗口计数和答案。若左端发生回退，或者窗口失去合法性后仍更新答案，就能很快定位错误。

\textbf{复杂度变化。} 暴力 O(nk)，固定窗口维护 O(n*26)，因字符集固定通常记作 O(n)，空间 O(26)。

\textbf{迁移追问。} 如果题目改成“维护差异计数可以避免每次比较整个数组”，先判断原状态是否仍然覆盖新答案集合，再决定是微调边界、改 value 语义，还是换算法。

\noindent\textbf{LeetCode 567 字符串的排列。}

\textbf{题目说明。} LeetCode 567 字符串的排列 的题面入口要先还原成具体任务：目标串任意排列出现在源串中，等价于固定长度窗口频次相同。

\textbf{样例入口。} 拿 s=eidbaooo、p=ab 手算，窗口 ba 的频次和 ab 相同，返回真；窗口长度固定为 2，每步只进一个字符出一个字符。

\textbf{暴力基线。} 暴力枚举 s2 中每个长度为 s1.size() 的窗口，重新统计频次并与 s1 频次比较；或者枚举 s1 全排列再查找。

\textbf{暴力过程。} 以 s1=ab、s2=eidbaooo 为例，固定长度窗口为 ei、id、db、ba、ao、oo、oo，其中 ba 命中。

\textbf{重复工作。} 题面模型：目标串任意排列出现在源串中，等价于固定长度窗口频次相同。暴力版的慢点要落到本题：浪费和 438 一样：相邻固定窗口只差一个进来字符和一个出去字符，不需要重新统计整个窗口，更不能枚举排列。后面的优化只消掉这类重复工作，不能改变答案集合。

\textbf{优化条件。} 本题的关键观察是：窗口长度固定为目标长度，右进左出维护字符计数差异。这句话必须能翻译成可维护状态；如果题目条件改变后观察不再成立，就不能继续套原来的写法。

\textbf{相邻状态。} 规律是排列出现等价于固定长度频次相同，不需要枚举所有排列。把这个特例继续拆开看：右端进入只带来一个新元素，左端离开只删掉一个旧元素；只有当旧左端被移走后不可能再产生更优答案时，窗口才可以单向推进。若这个单调淘汰条件不存在，就要换成前缀和、队列或其他状态。

\textbf{状态复用。} k=s1.size()。维护 26 长度频次数组或差异数组，窗口每右移一步执行加右减左；频次相同即可返回 true。手算时只改被新输入影响的那一小部分，而不是回头重算完整候选。

\textbf{指针和字段。} 状态字段要能说清：right 是新进入字符位置，right-k 是离开字符位置，diff/计数数组表示当前窗口和目标频次的差异。手算时按这个协议跟踪：窗口 db 右移到 ba 时，d 出去、a 进来；此时 b 和 a 各一次，与 ab 频次一致，返回 true。

\textbf{代码落点。} 关键代码是 for r from k to n-1: add s2[r], remove s2[r-k]。若用 matches 计数，可避免每轮比较 26 项。关键代码行必须能逐句翻译成“旧状态怎样变成新状态”，否则就只是模板。

\textbf{正确性。} s1 的某个排列出现在 s2 中，当且仅当存在长度 k 的窗口频次与 s1 完全相同；固定窗口维护精确覆盖所有这类窗口。

\textbf{边界实验。} 先用目标比源串长、重复字符、字符集固定、窗口重叠做反例检查。实验时重点跑目标比源串长、重复字符、字符集固定、窗口重叠，并打印左右端、窗口计数和答案。若左端发生回退，或者窗口失去合法性后仍更新答案，就能很快定位错误。

\textbf{复杂度变化。} 暴力 O(nk) 或排列阶乘级；固定窗口 O(n*26) 约 O(n)，空间 O(26)。

\textbf{迁移追问。} 如果题目改成“与找到所有异位词同型，只是返回布尔值”，先判断原状态是否仍然覆盖新答案集合，再决定是微调边界、改 value 语义，还是换算法。

\noindent\textbf{LeetCode 713 乘积小于 K 的子数组。}

\textbf{题目说明。} LeetCode 713 乘积小于 K 的子数组 的题面入口要先还原成具体任务：正整数数组的窗口乘积随右端扩展变大、随左端收缩变小。

\textbf{样例入口。} 拿 nums=[10, 5, 2, 6]、k=100 手算，右端到 2 时乘积 100 不满足严格小于，左端移走 10 后窗口 [5, 2] 合法，以 2 结尾的合法子数组有 [2] 和 [5, 2] 两个。

\textbf{暴力基线。} 暴力枚举每个起点 left，向右连续乘积 product，product<k 就计数，否则停止或继续检查。

\textbf{暴力过程。} 以 [10, 5, 2, 6]、k=100 为例，从 10 开始可数 [10]、[10, 5]，到 [10, 5, 2] 乘积为 100 不满足；从 5 又重新乘 5、5*2、5*2*6。

\textbf{重复工作。} 题面模型：正整数数组的窗口乘积随右端扩展变大、随左端收缩变小。暴力版的慢点要落到本题：浪费在相邻起点重复计算乘积。数组是正整数，乘入右端会让乘积不减，除出左端会让乘积不增。后面的优化只消掉这类重复工作，不能改变答案集合。

\textbf{优化条件。} 本题的关键观察是：右端乘入新数，乘积不小于 k 时左端不断除出，合法后以 right 结尾的子数组数为 right-left+1。这句话必须能翻译成可维护状态；如果题目条件改变后观察不再成立，就不能继续套原来的写法。

\textbf{相邻状态。} 规律是正整数乘积窗口具有单调性，每个右端贡献 right-left+1 个合法起点。把这个特例继续拆开看：右端进入只带来一个新元素，左端离开只删掉一个旧元素；只有当旧左端被移走后不可能再产生更优答案时，窗口才可以单向推进。若这个单调淘汰条件不存在，就要换成前缀和、队列或其他状态。

\textbf{状态复用。} left/right 维护乘积 product。每次 product*=nums[right]；当 product>=k 时不断 product/=nums[left++]。合法后新增 right-left+1 个以 right 结尾的子数组。手算时只改被新输入影响的那一小部分，而不是回头重算完整候选。

\textbf{指针和字段。} 状态字段要能说清：product 是当前窗口乘积，left 是最小合法起点，right-left+1 是以 right 结尾的合法起点数量。手算时按这个协议跟踪：到 right=2 时 product=100，不满足，移走 10 得 product=10，窗口 [5, 2] 合法，新增 [2] 和 [5, 2] 两个。

\textbf{代码落点。} 关键是 k<=1 时直接返回 0，因为所有 nums 为正，乘积不可能严格小于 1。计数用 ans += right-left+1。关键代码行必须能逐句翻译成“旧状态怎样变成新状态”，否则就只是模板。

\textbf{正确性。} 在 product<k 时，当前窗口的任意后缀乘积也小于 k，所以以 right 结尾且起点在 [left, right] 的子数组全合法。

\textbf{边界实验。} 先用k 小于等于一、元素都为一、乘积溢出、窗口缩到空、严格小于而不是小于等于做反例检查。实验时重点跑k 小于等于一、元素都为一、乘积溢出、窗口缩到空、严格小于而不是小于等于，并打印左右端、窗口计数和答案。若左端发生回退，或者窗口失去合法性后仍更新答案，就能很快定位错误。

\textbf{复杂度变化。} 暴力 O(\ensuremath{n^{2}})，滑动窗口 O(n) 时间、O(1) 空间。

\textbf{迁移追问。} 如果题目改成“若数组含零或负数，乘积单调性被破坏，需要换模型”，先判断原状态是否仍然覆盖新答案集合，再决定是微调边界、改 value 语义，还是换算法。

\noindent\textbf{LeetCode 862 和至少为 K 的最短子数组。}

\textbf{题目说明。} LeetCode 862 和至少为 K 的最短子数组 的题面入口要先还原成具体任务：数组允许负数，普通滑动窗口的和不再随左端移动单调变化。

\textbf{样例入口。} 拿 [2, -1, 2]、K=3 手算，普通正数窗口会被 -1 破坏单调性；前缀和为 0, 2, 1, 3，只有 3-0 达到 3。

\textbf{暴力基线。} 暴力枚举所有 left/right，用前缀和 O(1) 求区间和，检查是否至少为 K 并更新最短长度。

\textbf{暴力过程。} 以 [2, -1, 2]、K=3 为例，前缀和是 [0, 2, 1, 3]。普通正数窗口会被 -1 破坏：右端增加后和可能变小。

\textbf{重复工作。} 题面模型：数组允许负数，普通滑动窗口的和不再随左端移动单调变化。暴力版的慢点要落到本题：浪费在每个右端都回头扫描所有历史左端。需要的是某些前缀和尽量小且下标尽量靠后的候选左端。后面的优化只消掉这类重复工作，不能改变答案集合。

\textbf{优化条件。} 本题的关键观察是：在前缀和数组上用单调队列维护可能成为最优左端的下标。这句话必须能翻译成可维护状态；如果题目条件改变后观察不再成立，就不能继续套原来的写法。

\textbf{相邻状态。} 规律是在前缀和上用单调队列维护可能的最优左端，队首能满足时结算长度。把这个特例继续拆开看：右端进入只带来一个新元素，左端离开只删掉一个旧元素；只有当旧左端被移走后不可能再产生更优答案时，窗口才可以单向推进。若这个单调淘汰条件不存在，就要换成前缀和、队列或其他状态。

\textbf{状态复用。} 构造前缀和 prefix。用单调递增队列保存可能作为左端的前缀下标；当前前缀减队首前缀 >=K 时更新答案并弹出队首；队尾前缀不小于当前前缀时被当前支配，弹出。手算时只改被新输入影响的那一小部分，而不是回头重算完整候选。

\textbf{指针和字段。} 状态字段要能说清：deque 中保存前缀下标，且对应 prefix 值递增；i 是当前右端后一位，i-deque.front 是候选长度。手算时按这个协议跟踪：prefix 到 3 时，3-0>=3，长度 3；若出现更小的新前缀，它会支配队尾较大前缀，因为未来用它能得到更大区间和且更短或不更差。

\textbf{代码落点。} 关键有两个 while：先从队首结算满足 K 的最短长度，再从队尾删除 prefix 值不小于当前值的被支配候选。关键代码行必须能逐句翻译成“旧状态怎样变成新状态”，否则就只是模板。

\textbf{正确性。} 队首是当前最早且前缀较小的候选；一旦满足 K，继续保留它只会让未来长度更长，所以可以弹出。队尾被更小且更新的前缀支配，永远不会更优。

\textbf{边界实验。} 先用负数、答案不存在、K 为负或零、队尾支配关系做反例检查。实验时重点跑负数、答案不存在、K 为负或零、队尾支配关系，并打印左右端、窗口计数和答案。若左端发生回退，或者窗口失去合法性后仍更新答案，就能很快定位错误。

\textbf{复杂度变化。} 暴力 O(\ensuremath{n^{2}})，单调队列中每个前缀入队出队各一次，时间 O(n)，空间 O(n)。

\textbf{迁移追问。} 如果题目改成“这是从窗口题升级到前缀和加单调队列的代表”，先判断原状态是否仍然覆盖新答案集合，再决定是微调边界、改 value 语义，还是换算法。

\noindent\textbf{LeetCode 1438 绝对差不超过限制的最长连续子数组。}

\textbf{题目说明。} LeetCode 1438 绝对差不超过限制的最长连续子数组 的题面入口要先还原成具体任务：窗口合法性取决于窗口内最大值和最小值之差。

\textbf{样例入口。} 拿 [8, 2, 4, 7]、limit=4 手算，窗口 [8, 2] 最大最小差 6，不合法，左端必须右移；窗口 [2, 4] 合法。

\textbf{暴力基线。} 暴力枚举每个窗口，重新扫描窗口求最大值和最小值，若 max-min<=limit 就更新最长长度。

\textbf{暴力过程。} 以 [8, 2, 4, 7]、limit=4 为例，窗口 [8, 2] 差为 6 不合法，窗口 [2, 4] 合法。暴力每次都重新找最大最小。

\textbf{重复工作。} 题面模型：窗口合法性取决于窗口内最大值和最小值之差。暴力版的慢点要落到本题：浪费在相邻窗口重复求极值。右端新增一个元素后，只有被它支配的旧极值候选需要删除；左端移动时只需清理过期下标。后面的优化只消掉这类重复工作，不能改变答案集合。

\textbf{优化条件。} 本题的关键观察是：两个单调队列分别维护最大和最小，超过限制时移动左端并清理过期下标。这句话必须能翻译成可维护状态；如果题目条件改变后观察不再成立，就不能继续套原来的写法。

\textbf{相邻状态。} 规律是两个单调队列分别维护最大值和最小值，差值超限时移动左端并清理过期下标。把这个特例继续拆开看：右端进入只带来一个新元素，左端离开只删掉一个旧元素；只有当旧左端被移走后不可能再产生更优答案时，窗口才可以单向推进。若这个单调淘汰条件不存在，就要换成前缀和、队列或其他状态。

\textbf{状态复用。} 用 maxDeque 维护递减下标队列，队首是窗口最大值；minDeque 维护递增下标队列，队首是窗口最小值。差值超限时移动 left 并清理过期。手算时只改被新输入影响的那一小部分，而不是回头重算完整候选。

\textbf{指针和字段。} 状态字段要能说清：left/right 是窗口边界，maxDeque.front/minDeque.front 是当前极值下标，answer 是最长合法窗口长度。手算时按这个协议跟踪：加入 2 后，最大队首为 8、最小队首为 2，差 6 超限；left 从 0 移到 1，过期的 8 被清理，窗口 [2] 合法。

\textbf{代码落点。} 关键是入队前从队尾弹出被新值支配的下标；收缩时如果队首下标小于 left 就弹出。关键代码行必须能逐句翻译成“旧状态怎样变成新状态”，否则就只是模板。

\textbf{正确性。} 单调队列保留的都是可能成为当前或未来极值的下标，被新元素支配的旧下标既更旧又不更优，永远不会再成为极值。

\textbf{边界实验。} 先用limit 为零、重复元素、最大最小同时过期、窗口收缩多次做反例检查。实验时重点跑limit 为零、重复元素、最大最小同时过期、窗口收缩多次，并打印左右端、窗口计数和答案。若左端发生回退，或者窗口失去合法性后仍更新答案，就能很快定位错误。

\textbf{复杂度变化。} 暴力 O(\ensuremath{n^{2}})，两个单调队列每个下标入队出队各一次，时间 O(n)，空间 O(n)。

\textbf{迁移追问。} 如果题目改成“它说明窗口状态可以是动态极值，不只是频次或总和”，先判断原状态是否仍然覆盖新答案集合，再决定是微调边界、改 value 语义，还是换算法。

\subsection{实验复盘}

追问通常会问窗口为什么不会漏掉左端、为什么有负数会失效、固定窗口和可变窗口有什么区别。请准备一个反例证明错误窗口条件会漏解。

复盘本节时先从 LeetCode 3 无重复字符的最长子串、LeetCode 11 盛最多水的容器、LeetCode 15 三数之和 开始：每道题都先手写一个能覆盖全部候选的暴力版，再指出优化结构具体保存了哪一段历史信息，最后用相邻案例比较为什么不能互换解法。
